% ======================================================================
%  21_layer1_geom.tex — 结构层（下）：§6–§12 的结构卡
%  只写正文，preamble 由 main.tex 提供；严格按 plan/01 §5 六栏卡格式。
%  所有锚点页码均为印刷页码 N（正文在 source/text/p{N+1:03d}.txt）。
% ======================================================================

% 本节局部记号（preamble 不改）：对偶群的权与表示记为 ^L w、^L R。

\section{结构层（下）：从膜算子到 $D$-模翻译器（\S6--\S12）}

\subsection{\S6--\S12 总述：一条四段链、三种可信度}

结构层上半（\S2--\S5）拆的是机器：如何把 $N=4$ 超对称 Yang--Mills 扭化成 4d 拓扑场论，如何在
Riemann 面 $C$ 上紧化得到 $\MH$ 上的 2d 拓扑 sigma 模型，以及这个模型里"膜"是什么意思。下半
（\S6--\S12）拆的是用法：把这台机器对准几何 Langlands 的两个方向。与上半不同，这七节不是七个并列
话题，而是一条递进链；每一节的输出是下一节的输入，任何一环断掉，几何 Langlands 的标准陈述就拼不
起来。本节给这七节各画一张结构卡，并在每张卡上标出：哪些陈述是原文的硬结论（A），哪些是物理上自
然而原文一句话带过的推断（B），哪些是作者自己说 "we expect / we believe" 的猜想（C）。

\begin{keyidea}
\textbf{链条一句话}：\S6 把 4d 线算子变成 2d 膜范畴的自函子（骨架 S5）；\S7--\S8 证明电特征膜就是
同态 $\pi_1(C)\to\LdualC$，其中 \S7 的 flux 修正不可省（S6）；\S9--\S10 证明 't Hooft 算子的作用
恰好是 Hecke 修改（S7）；\S11--\S12 用典范余各向同性 A-膜把 magnetic 侧翻译成 $D$-模，得到几何
Langlands 的标准形式（S8）。
\end{keyidea}

\textbf{(1) S5，\S6：线算子 $=$ 膜算子。}这一步的真正发明不是"线算子存在"（任何规范理论都有 Wilson
线），而是\textbf{把线算子当成作用在膜范畴上的算符}，并因此只能说"特征值是向量空间"：$XB\simeq B\otimes V$
\pA\anchor{src p.95}。没有这个升级，"本征值"就只是数，几何 Langlands 里的 eigensheaf（特征层）在物理
侧没有对应物。4d 来源还白送两个结构：线算子自动交换，因此可以同时对角化；特征值 $V_p$ 随 $p\in C$
成为平坦丛，已经预演了"局部系统"。

\textbf{(2) S6，\S7--\S8：电特征膜 $=\pi_1(C)$ 的表示。}物理侧 \S8.2 只用一行计算——任意向量丛在
一点处平凡——就得到：位于 $x\in\MH$ 的 zerobrane 是电特征膜，特征值是 $E_q(R)|_x$ \pA\anchor{src p.121}。
但 \S7 必须插在中间，因为 universal $G$-bundle 不存在，Chan-Paton 丛只能是 twisted 的，离散电
flux $e_0$ 恰好就是 twist 的标签；正确的物理对象是"所有 $e_0$ 的 twisted zerobrane 之和" \pA\anchor{src p.122}。
删掉 \S7，\S8.2 的"特征值"在 $e_0\ne0$ 时根本没有意义。这是七节链里最容易被数学读者误判为"背景复
习"的一节。

\textbf{(3) S7，\S9--\S10：'t Hooft $=$ Hecke。}\S9.1 对 $U(N)$ 把 't Hooft 算子的奇点逐分量解出来，
得到 $L_i\mapsto L_i\otimes\mathcal O(p_0)^{m_i}$，这就是 Hecke 修改的定义式 \pA\anchor{src p.132}。
\S9.2--\S9.3 说明"所有可能的 Hecke 修改"构成仿射 Grassmannian 的变体；\S10 用 Bogomolny 方程把这
个空间实现成边界条件的模空间，并在 \S10.4 算出 't Hooft 算子的 OPE。这一步补上了磁侧缺失的表示论：
Wilson 有现成的 $\mathrm{Rep}(G)$，而 't Hooft 的代数结构必须由 Hecke 修改空间的几何给出。

\textbf{(4) S8，\S11--\S12：$D$-模翻译器。}\S11.1 在 $\MH$ 上造出典范余各向同性 A-膜，\S11.2 说明
A-膜给出 $\MC$ 上的 $D$-模；于是 \S8.2 的 type F 膜（magnetic eigenbrane）对应 $D$-模，正是几何
Langlands 的标准陈述。\S12 不是新定理：它把全篇的膜放进一张按三个复结构分类的检查清单，供读者核对
每种膜落在 GL 的哪一边。因此 \S12 的卡片要点不是"得到结论"，而是"如何不重不漏地列举膜"。

\paragraph{与 \S2--\S5 的文本差异：数学读者必须换一种读法。}
上半的 A 级陈述（扭化、超凯勒商、Hitchin 纤维化的 T-对偶）基本可以当定理读；下半的 A 级陈述要打
折，因为整条链建立在两个物理前提上：(i) $N=4$ 超对称 Yang--Mills 的 S-对偶本身是一个有强证据但未
证明的猜想，原文在 \S2.2 把它当作已建立的事实使用 \pA\anchor{src p.11}；(ii) 4d 拓扑场论的存在与
拓扑性来自 \S3 的物理构造，其中 \S3.3 的消失定理只在附加条件下成立 \pA\anchor{src p.28}。所以本书
标 A 的准确含义是"原文明确写出、且我们核验了其推理"，不是"数学上已被证明"。下半真正的"猜点"集中在
以下几处，读到那里应放慢：

\begin{table}[htbp]
\centering\small
\begin{tabular}{@{}p{1.85cm}p{5.1cm}p{8.45cm}@{}}
\toprule
\tabhead{位置} & \tabhead{原文语气} & \tabhead{读者要小心什么} \\
\midrule
\S6.4 & "we expect"（(6.28)--(6.29)） & 't Hooft 算子的 OPE 没有物理侧的直接计算，作者只从 Coulomb 分支读出最高权部分，其余交给 \S10.4 的 Hecke 修改空间几何 \pC\anchor{src p.99} \\
\S9.2 & "we will see / it is natural" & Hecke 修改空间的连通性与维数依赖单极子解的存在性；数学上这是已知定理，但物理侧是重建 \pB\anchor{src p.135} \\
\S10.2 & "monopole bubbling" & 单极子成泡的图像是半经典推断，原文用它解释 Hecke 空间的奇点 \pC\anchor{src p.155} \\
\S11.2 & "we argue / we believe" & A-膜 $\leftrightarrow$ $D$-模的对应是启发式论证，不是范畴等价的证明 \pC\anchor{src p.196} \\
\S12 & "we expect / we believe" & 三种膜型的清单是物理分类，其完备性没有独立证明 \pC\anchor{src p.207} \\
\bottomrule
\end{tabular}
\caption{后半程的"猜点"清单。标 C 的断言在正文卡片里都会再次点明，不要把它们当定理引用。}
\label{tab:geom-guesses}
\end{table}

\begin{figure}[htbp]
\centering
\begin{tikzpicture}[node distance=5mm and 5mm, box/.style={draw=kwblue!65, fill=kwblue!4,
  rounded corners=2pt, align=center, font=\small, inner sep=4pt, text width=3.25cm}]
\node[box] (s5) {\textbf{S5}（\S6）\\ 线算子 $=$ 膜算子\\[2pt]
  {\scriptsize 输出：自函子 $X$、$XB=B\otimes V$}};
\node[box, right=of s5] (s6) {\textbf{S6}（\S7--\S8）\\ 电特征膜\\[2pt]
  {\scriptsize 输出：$\pi_1(C)\to\LdualC$ 的表示}};
\node[box, right=of s6] (s7) {\textbf{S7}（\S9--\S10）\\ 't Hooft $=$ Hecke\\[2pt]
  {\scriptsize 输出：Hecke 算子代数}};
\node[box, right=of s7] (s8) {\textbf{S8}（\S11--\S12）\\ c.c.\ A-膜 $\to$ $D$-模\\[2pt]
  {\scriptsize 输出：GL 标准形式}};
\draw[-{Latex[length=2mm]}] (s5) -- node[above,font=\scriptsize]{Wilson} (s6);
\draw[-{Latex[length=2mm]}] (s6) -- node[above,font=\scriptsize]{对偶} (s7);
\draw[-{Latex[length=2mm]}] (s7) -- node[above,font=\scriptsize]{magnetic} (s8);
\end{tikzpicture}
\caption{后半程的四段链。箭头上的词是"上一段的输出被下一段消费的部分"；每段内部还可细分：\S7 是
\S8 的修正项，\S10 是 \S9 的几何实现，\S12 是 \S11 的分类清单。}
\label{fig:geom-chain}
\end{figure}

% ======================================================================
\subsection{\S6 圈与线算子：从 4d 算子到 2d 膜范畴（pp.78--102）}

\S6 是后半程的入口，也是全文语气转折的地方：从这里开始作者不再造理论，而是让理论去作用。本节的
输出是骨架 S5——4d 拓扑场论的线算子给出 2d 膜范畴的自函子，特征对象（eigenbrane）的"特征值"是一个
向量空间。这一步的输入是 \S3 的标量超荷与扭化参数、\S4 的 $\MH$ 三复结构、\S5 的膜范畴；出口分三
路：\S7--\S8 修正并具体化电特征膜，\S9--\S10 具体化磁侧的 Hecke 作用，\S11 要用的"magnetic
eigenbrane 带平坦丛"也在本节末尾就位。

\S6 内部四小节的逻辑是：先造两种基本算子（Wilson 在 $\Psi=\infty$，'t Hooft 在 $\Psi=0$，混合算子
在有理 $\Psi$），再把它们紧化到 2d 并翻译成"作用在膜上"，最后定义 eigenbrane 并证明 4d 来源给出的
两个自动性质。注意 6.1 与 6.2 的对称性：Wilson 的做法是"把 $\varphi$ 加进联络"，'t Hooft 的做法是
"让 $\varphi$ 跟着奇点发散"；两者都是为了让同一个标量超荷 $Q$ 存活。

\begin{kcard}{\S6.1 拓扑 Wilson 算子（pp.79--80）}
\begin{enumerate}[label=\textbf{\arabic*}]
\item \textbf{这一节的任务}：在 $\Psi=\infty$（即 $t=\pm\ii$）构造保持拓扑超荷 $Q$ 的 Wilson 线算子，
  并说明为什么普通的 $W_0(R,S)=\Tr_R P\exp(-\oint_S A)$ 不够、必须把规范场替换成
  $\mathcal A=A+\ii\varphi$。
\item \textbf{输入}：\S3.1 的拓扑变换 (6.2)--(6.3)；骨架 S2（存在标量超荷 $Q$）、S3（$M=\Sigma\times C$
  与 $\MH$ 的几何）。
\item \textbf{关键构造}：普通 Wilson 圈在拓扑变换 $A_\mu\mapsto A_\mu+\ii u\psi_\mu+\ii v\tilde\psi_\mu$
  下不变；但 Higgs 场按 $\delta_T\varphi_\mu=\ii v\psi_\mu-\ii u\tilde\psi_\mu$ 变换，两者在
  $t=v/u=\pm\ii$ 时的组合不变：
\begin{equation}
\begin{aligned}
W(R,S)&=\Tr_R P\exp\Big(-\oint_S (A+\ii\varphi)\Big) &&(t=\ii),\\
W(R,S)&=\Tr_R P\exp\Big(-\oint_S (A-\ii\varphi)\Big) &&(t=-\ii).
\end{aligned}
\label{eq:geom-wilson}
\end{equation}
  开线情形的 $W(R,S)$ 不是迹，而是从 $p$ 到 $q$ 的平行移动矩阵 \pA\anchor{src p.79}\anchor{src p.80}。
\item \textbf{结论}：Wilson 算子是 $\Psi=\infty$ 处保持拓扑对称的线算子，按 $G$ 的不可约表示 $R$ 分类；
  其存在性的一句话解释是：$t=\pm\ii$ 的拓扑方程断言 $\mathcal A$ 平展，故沿闭圈的 holonomy 是自然
  不变量 \pA\anchor{src p.80}。
\item \textbf{出口}：\S8.1 要把这个算子对膜的作用算成 Chan-Paton 丛的张量积；为此所需的"$\varphi$ 的
  费米补偿"正是 \eqref{eq:geom-wilson} 里 $\ii\varphi$ 的 2d 对应物 \pB\anchor{src p.119}。
\item \textbf{锚点索引}：p.79 定义 $W_0(R,S)$ 与符号约定；p.80 变换律 (6.2)--(6.3) 与 $t=\pm\ii$ 的
  组合；p.80 开线算子 $=$ 平行移动矩阵；p.80 拓扑方程 $\Rightarrow$ 平展性。
\end{enumerate}
\textbf{可信度}：公式与变换律是 A。"在 $t=\pm\ii$ 存在极自然"是作者的物理判断（B），依据只有平展性；
"对一般 $t$ 不存在"作者自己加了限定——\S3.3 的消失定理需要全局输入、对有膜的场合不一定成立
\pB\anchor{src p.80}。
\end{kcard}

\begin{kcard}{\S6.2 拓扑 't Hooft 算子（pp.81--85）}
\begin{enumerate}[label=\textbf{\arabic*}]
\item \textbf{这一节的任务}：构造 Wilson 算子在 S-对偶下的对象——'t Hooft 算子，并确定它在 $\Psi=0$
  保持 $Q$ 所需的场奇点。这是"序算子 $\leftrightarrow$ 无序算子"在拓扑理论中的实现。
\item \textbf{输入}：\S2.2 的 S-对偶 $S:\Psi\mapsto-1/n_g\Psi$；\S6.1 的 Wilson 算子；GNO 对偶
  （`\S\S1` 想法 3 的算子版）。
\item \textbf{关键构造}：$U(1)$ 情形取 Dirac 单极奇点
\begin{equation}
F=\frac{\ii}{2}\,\star_3 \dd\frac{1}{r},\qquad
\varphi=\frac{\ii}{2|\vec x|}\,\dd s ,
\label{eq:geom-thooft}
\end{equation}
  其中 $\star_3$ 是法向 $\RR^3$ 的 Hodge 星、$\dd s$ 是沿线方向的标架；带电荷 $m$ 时整体系数乘 $m$。
  一般 $G$ 用同态 $\rho:U(1)\to G$（共轭类 $=$ $\Ldual$ 的最高权 $=$ $\Ldual$ 的不可约表示）嵌入
  \eqref{eq:geom-thooft}，得到算子 $T(\LR,S)$ \pA\anchor{src p.83}\anchor{src p.84}。
\item \textbf{结论}：'t Hooft 算子按 $\Ldual$ 的不可约表示分类（GNO 对偶的算子版）；其奇点携带特征
  类 $\xi(\LR)\in H^2(S^2,\pi_1(G))$，拓扑非平凡当且仅当 $\Ldual$ 的中心在 $\LR$ 上非
  平凡地作用 \pA\anchor{src p.84}。
\item \textbf{出口}：\S9.1 把这个奇点逐分量解出来，发现它的效果恰好是 Hecke 修改；\S8.1 用 $\xi(\LR)$
  写出 't Hooft 算子如何改变 $m_0$ \pA\anchor{src p.120}。
\item \textbf{锚点索引}：p.81 't Hooft 作为 disorder 算子的定义（法丛定向）；p.82--p.83 单极奇点与
  Maxwell 方程；p.83 (6.9) 保持拓扑对称的 $\varphi$ 奇点；p.84 $\rho:U(1)\to G$、GNO 分类、与
  $\PP^1$ 上 $G$-丛分类的关系。
\end{enumerate}
\textbf{可信度}：奇点公式与 GNO 分类是 A（引 [10],[90]）。"$\PP^1$ 上 $G$-丛的 Grothendieck 分类
解释 't Hooft 算子"是 B：论证依赖"$\PP^1$ 上 Yang--Mills 解都等价于阿贝尔解"这一物理陈述
\pB\anchor{src p.84}。而"$\Psi=0$ 必有对偶算子"本身依赖 S-对偶（见总述），属于条件性的 A。
\end{kcard}

\begin{kcard}{\S6.3 紧化到二维（pp.89--92）}
\begin{enumerate}[label=\textbf{\arabic*}]
\item \textbf{这一节的任务}：确定 4d 线算子紧化到 $\Sigma$ 之后在 2d 有效理论里的身份——答案是把线
  算子翻译成"边界条件/膜"，这是全篇"算子作用于膜"的语言基础。
\item \textbf{输入}：S3（$M=\Sigma\times C$ 紧化）；\S6.1--\S6.2 的两种算子；\S5 的膜与 $(B_1,B_2)$
  开弦空间。
\item \textbf{关键构造}：按曲线 $S\subset\Sigma\times C$ 的形状分三类：(a) $S=p\times S'$（在 $C$ 内绕
  圈）在 2d 中是点状局部算子；(b) $S=S''\times q$（沿 $\Sigma$ 走）仍是线算子；(c) 一般曲线同伦于
  "在 $\Sigma$ 上走、间或在 $C$ 中绕圈"，故是线算子加上插在特殊点的局部算子 \pA\anchor{src p.89}\anchor{src p.90}。
  关键的折叠（folding）技巧：一条线算子把平面分成两个区域 $Y$ 与 $Z$，把 $Z$ 的定向翻转后折叠，
  线算子变成 $Y\otimes Z$ 中的边界条件，即 $(Y,Z)$-膜 \pA\anchor{src p.90}。由此引出
  $H_{B_1,B_2}$（$(B_1,B_2)$ 开弦空间）、拼接乘法 $H_{B_1,B_2}\otimes H_{B_2,B_3}\to H_{B_1,B_3}$，
  以及 $H_{B,B}$ 的结合代数结构 \pA\anchor{src p.91}。
\item \textbf{结论}：在 2d 中"线算子 $=$ 作用在膜上的算符"有了严格含义：$X$ 把 $B$ 变成新膜 $XB$；
  边界上的局部算子空间就是 $H_{B,B}$，而 $H_{B,\tilde B}$ 是 $(H_{B,B},H_{\tilde B,\tilde B})$ 双模
  \pA\anchor{src p.92}。
\item \textbf{出口}：\S6.4 的 eigenbrane 定义只需要这一节的"$X$ 可以作用在 $B$ 上"；\S8.1 则反过来
  用 2d 边界项的 holonomy 与 4d Wilson 线比较。
\item \textbf{锚点索引}：p.89 图 4 的三种曲线；p.90 折叠技巧与 $(Y,Z)$-膜；p.91 图 6 的拼接与结合律；
  p.92 边界局部算子 $=H_{B,B}$、双模结构。
\end{enumerate}
\textbf{可信度}：几何分解与折叠技巧是 A（原文写得完整，技巧引 [93--95]）。"$H_{B,B}$ 是结合代数"
是 TQFT 公理的直接应用，数学上是标准的（A）；但"边界局部算子空间 $=H_{B,B}$"的切割论证是物理图
像（B，见图 6(d) 的解释）\pB\anchor{src p.92}。
\end{kcard}

\begin{figure}[htbp]
\centering
\begin{tikzpicture}[font=\small,
  box/.style={draw=kwblue!65, fill=kwblue!4, rounded corners=2pt, align=center,
    inner sep=4pt, text width=3.4cm}]
\node[box] (op) {4d：线算子 $X$\\ 支撑在 $S\subset M$};
\node[box, right=1.2cm of op] (br) {2d：边界条件\\（膜 $B$）};
\node[box, right=1.2cm of br] (nb) {新膜 $XB$\\（$X$ 贴近边界）};
\draw[-{Latex[length=2mm]}] (op) -- node[above,font=\scriptsize]{紧化} (br);
\draw[-{Latex[length=2mm]}] (br) -- node[above,font=\scriptsize]{作用} (nb);
\node[draw=kwred!60, rounded corners=2pt, inner sep=4pt, font=\scriptsize,
  below=9mm of br, text width=6.2cm] (eig) {特征膜：$XB\simeq B\otimes V$，
  {\normalfont 特征值不是数而是向量空间 $V$}};
\draw[-{Latex[length=1.6mm]}, kwred!60, dashed] (nb.south) to[out=-90,in=0] (eig.east);
\end{tikzpicture}
\caption{线算子的两级翻译：先把 4d 线算子紧化成 2d 边界条件，再让它在膜上作用。第二级翻译是 \S6
的中心思想，也是"特征值升级为向量空间"的来源。}
\label{fig:geom-op-brane}
\end{figure}

\begin{kcard}{\S6.4 边界附近的线算子与 eigenbrane（pp.93--102）}
\begin{enumerate}[label=\textbf{\arabic*}]
\item \textbf{这一节的任务}：定义线算子作用在膜上（$B\mapsto XB$）与 eigenbrane（$XB=B\otimes V$），
  并证明"来自 4d"给 eigenbrane 带来两个自动结构：算子族自动交换、特征值随 $p\in C$ 构成平坦丛。
\item \textbf{输入}：\S6.3 的膜作用；\S5 的膜范畴；S3（线算子来自 4d，故可定义在任意 $p\in C$ 处）。
\item \textbf{关键构造}：线算子 $X$ 贴近边界时把边界条件 $B$ 变成有效边界条件 $XB$（p.94）；
  张量积 $B\otimes V$ 由 $H_{\tilde B,B\otimes V}=H_{\tilde B,B}\otimes V$ 刻画；于是
\begin{equation}
XB\simeq B\otimes V \qquad\text{（eigenbrane 的定义，对应 Beilinson--Drinfeld 的 eigensheaf）}.
\label{eq:eigenbrane}
\end{equation}
  两个算子交换的定义是 $X_1X_2B=X_2X_1B$（对一切 $B$），交换则可同时对角化；来自 4d 的线算子自动
  交换（论证见原文脚注 31：两算子的支撑点可在 $C$ 中分开、不产生奇点）。特征值随 $p$ 变化：
  $X_pB=B\otimes V_p$，且移动 $p$ 等价于改变 $C$ 的度量、只差 BRST 平凡项，故 $V_p$ 上有一个平坦
  联络——$V_p$ 是 $C$ 上平坦丛的纤维 \pA\anchor{src p.93}\anchor{src p.95}\anchor{src p.96}\anchor{src p.97}。
\item \textbf{结论}：\textbf{electric eigenbrane} 是 Wilson 算子族的联合特征膜，\textbf{magnetic
  eigenbrane} 是 't Hooft 算子族的；$S$ 交换二者，并交换两种膜的超对称型；(B,B,B) 与 (B,A,A) 的
  Wilson/'t Hooft 特殊线算子分别对应（见下卡）\pA\anchor{src p.97}\anchor{src p.98}。Wilson 算子族的
  OPE 就是 $G$ 的表示环（$R\otimes\tilde R=\oplus_\alpha n_\alpha R_\alpha$ 给出逐项乘积）；
  't Hooft 算子族的 OPE 与 $\Ldual$ 的表示环同构——\textbf{但这一步原文是猜的}
  \pC\anchor{src p.98}\anchor{src p.99}。
\item \textbf{出口}：\S8.2 把 zerobrane 具体地认成 electric eigenbrane；\S9--\S10 证明 't Hooft 族
  的作用 $=$ Hecke 算子族，从而把上式中的"交换代数"落实为 Hecke 算子代数。
\item \textbf{锚点索引}：p.94 $B'=XB$；p.95 (6.22) eigenbrane 定义；p.95--p.96 张量积与 eigensheaf；
  p.97 (6.25) 平坦丛与脚注 31 的交换性论证；p.98 (6.27) Wilson OPE；p.99 (6.29) 't Hooft OPE
  的期望；p.101 线算子 $=$ 膜范畴的自函子；p.102 special Wilson 的 (B,B,B) 型。
\end{enumerate}
\textbf{可信度}：定义与 Wilson OPE 是 A；"4d 来源 $\Rightarrow$ 自动交换"是 A$^-$——论证只在 $C$
至少二维时成立，作者自己指出这一点 \pB\anchor{src p.97}；'t Hooft OPE (6.29) 是 C：原文说
"we expect"，并明确承认物理文献没有直接计算过 't Hooft 算子的 OPE，只从 Coulomb 分支读出最高权
项，其余交给 \S10.4 \pC\anchor{src p.99}。special Wilson/'t Hooft 的 (B,B,B)/(B,A,A) 型是 A
（由 (6.30)--(6.31) 与 \S5.3 表 2 得到）\pA\anchor{src p.102}。
\end{kcard}

\begin{figure}[htbp]
\centering
\begin{tikzpicture}[font=\small, every node/.style={align=center}]
\node (I) at (90:2.05) {\textbf{复结构 $I$}};
\node (J) at (210:2.05) {\textbf{复结构 $J$}};
\node (K) at (330:2.05) {\textbf{复结构 $K$}};
\draw[kwblue!60, thick] (J) -- (I);
\draw[kwblue!60, thick] (I) -- (K);
\draw[kwblue!60, thick] (J) -- (K);
\node[font=\scriptsize, text width=4cm, above=2mm of I] {(B,B,A) 与 (A,A,A) 不存在};
\node[draw=kwblue!65, fill=kwblue!4, rounded corners=2pt, inner sep=3pt,
      text width=5.0cm, anchor=west, font=\scriptsize] at (3.6,0.1)
  {允许的四类词：\\
   (B,B,B)：0-膜、space-filling 膜\\
   (B,A,A)：Hitchin 纤维、$I$ 中复拉格朗日（$\Omega_I$ 全纯辛）\\
   (A,B,A)/(A,A,B)：$J/K$ 中复拉格朗日（$\omega_J,\omega_K$）};
\end{tikzpicture}
\caption{膜类型的三角关系：一个膜在三个复结构 $I,J,K$ 上分别取 B 型或 A 型，构成一个三位词。物理上
允许的是 (B,B,B)、(B,A,A)、(A,B,A)、(A,A,B)；若一个子流形在 $I,J$ 上全纯，则自动在 $K$ 上全纯，
故 (B,B,A) 不存在。special Wilson 线 $=(B,B,B)$，special 't Hooft 线 $=(B,A,A)$，$S$-对偶恰好交换
这两个词 \pA\anchor{src p.102}。}
\label{fig:geom-brane-types}
\end{figure}

% ======================================================================
\subsection{\S7 离散 flux 与 S-对偶：eigenbrane 的必要修正（pp.103--115）}

\S7 表面上是对离散电磁 flux 的"复习"，实际是 S6 的修正项：电特征膜的特征值是一个丛，但 universal
$G$-bundle 不存在的障碍使这个丛只能是 twisted 的，而 twist 的标签正是离散电 flux $e_0$。本节输出
两条断言：第一，universal bundle 作为 twisted vector bundle 存在，其 twist 由表示 $R$ 的中心特征
$\theta_R(\zeta)$ 给出；第二，\S7.2 把 4d 的 $(e_0,m_0,e_1,m_1)$ 逐项翻译成 2d sigma 模型的数据
（B-field、$\MH$ 的连通分支、守恒动量与绕数），从而给出 S-对偶在膜范畴上的精确形式。

\S7 的可信度分布值得注意：群论与拓扑部分（完美配对、Poincaré 对偶、twisted bundle）是 A；把 $e_0$
解释成 $\MH$ 上的平坦 B-field 是 B（作者给了逐项验证，见 p.111--p.112）；而 $e_1,m_1$ 的"守恒动量/
绕数"解释与最后关于两种量子化的评论是 C（作者用 "we can simply specify" 的语气给出选择）。

\begin{kcard}{\S7.1 离散 flux 与 universal bundle（pp.103--109）}
\begin{enumerate}[label=\textbf{\arabic*}]
\item \textbf{这一节的任务}：建立两个工具——中心 $Z$ 的自然完美配对（用于定义离散电磁 flux），以及
  universal $G$-bundle 的障碍类 $\zeta$（用于说明为什么 Chan-Paton 丛必须 twisted）。
\item \textbf{输入}：\S2.2 的 S-对偶群 $\mathrm{SL}(2,\ZZ)$；群论：$Z=\Lambda^\vee/\Lambda$；\S4 的
  $\MC$ 与 $\MH$。
\item \textbf{关键构造}：有限群 $Z$ 的自然自对偶配对
\begin{equation}
\Upsilon:Z\times Z\to U(1),\qquad \Upsilon(a,b)=\exp(2\pi\ii\langle a,b\rangle),
\label{eq:geom-pairing}
\end{equation}
  由此 $\hat\Upsilon$ 使 $H^2(M,\ZZ)$ 自对偶；$H^2(M,\ZZ)\ni\xi(E)$ 在 $M=S^1\times W$ 上分裂为
  $a\oplus m$，离散电、磁 flux 是 $(e,m)\in \mathrm{Hom}(H^1(W,\ZZ),U(1))\times H^2(W,\ZZ)$，
  S-对偶作用 $S:(e,m)\mapsto(-m,e)$ \pA\anchor{src p.104}\anchor{src p.105}\anchor{src p.106}。
  universal bundle：对手伴随群 $G_{\mathrm{ad}}$ 存在；对单连通 $G$，过渡函数在 triple overlap 上
  只差中心元素 $\Theta_{ki}\Theta_{jk}\Theta_{ij}=f_{ijk}\in Z$，给出障碍类
  $\zeta=\xi(E_{\mathrm{ad}})\in H^2(M,\ZZ)$；它作为 twisted vector bundle 总存在，twist 为
  $\varrho(\zeta)\in H^2(M,U(1))$ \pA\anchor{src p.107}\anchor{src p.108}\anchor{src p.109}。
\item \textbf{结论}：对非伴随型 $G$，\textbf{不存在}通常意义的 universal $G$-bundle，但存在 twisted
  的版本；twist 的标签由表示 $\varrho$ 的中心特征决定。这正是 \S8.1 中 Wilson 算子改变离散电场的
  数学载体。
\item \textbf{出口}：\S7.2 把 $e_0$ 解释为 B-field、把 $m_0$ 解释为 $\MH$ 的连通分支；\S8.1 用
  $E_p(R)$ 的 twist 写出 $e_0\to e_0+\theta_R$。
\item \textbf{锚点索引}：p.104 中心表与完美配对 (7.1)--(7.2)；p.105--p.106 离散 flux 与 $S:(e,m)\mapsto(-m,e)$；
  p.107--p.108 universal bundle 的拼接过障碍；p.109--p.110 twisted vector bundle 与 $\varrho(\zeta)$。
\end{enumerate}
\textbf{可信度}：配对与 flux 的 S-对偶变换是 A（标准拓扑与 [45,101]）。universal bundle 障碍的构造是
A（作者自陈是"pedestrian"的初等论述）\pA\anchor{src p.108}。"用 twisted bundle 解释 Chan-Paton"是 B：
这一步在 \S7.2--\S8.1 才被真正使用，原文在此只是铺垫 \pB\anchor{src p.112}。
\end{kcard}

\begin{kcard}{\S7.2 紧化到二维：flux 的 2d 含义（pp.110--115）}
\begin{enumerate}[label=\textbf{\arabic*}]
\item \textbf{这一节的任务}：把 4d 的离散 flux 逐项翻译成 2d sigma 模型（target $\MH$）的数据，并推出
  S-对偶在膜上的作用规律。这是 S6 中"flux 修正不可省"的正式版本。
\item \textbf{输入}：\S7.1；\S5.4--\S5.5（S-对偶 $=$ 镜像对称、Hitchin 纤维化）；\S4 的 $\MH$ 连通
  分支与 Hitchin 纤维。
\item \textbf{关键构造}：$\xi(E)\in H^2(\Sigma\times C,\ZZ)$ 按 Künneth 分解
\begin{equation}
\xi(E)=\xi_{2,0}+\xi_{1,1}+\xi_{0,2}\in H^2(\Sigma)\oplus H^1(\Sigma,H^1(C))\oplus H^2(C),
\label{eq:geom-kunneth}
\end{equation}
  并与低能 map $\Phi:\Sigma\to\MH$ 联系：$\xi(E)=(\Phi\times1)^*(\zeta)$ \pA\anchor{src p.110}。
  逐项翻译：$m_0=\xi_{0,2}$ 是 $\MH(G_{\mathrm{ad}})$ 的连通分支标签；$e_0$ 给出
  $\MH(G_{\mathrm{ad}})$ 上的平坦 B-field $b_{e_0}=e_0(\zeta)$，在存在 B-field 时 Chan-Paton 丛变成
  twisted 的 \pA\anchor{src p.111}\anchor{src p.112}。S-对偶交换 $(e_0,m_0)\mapsto(-m_0,e_0)$：
  一侧的"分支拓扑"对应另一侧的"B-field twist"；特例给出 Hitchin 纤维 $F$ 与 $\Ldual F$ 之间的
  twisted duality（Hausel--Thaddeus 定理）\pA\anchor{src p.113}\anchor{src p.114}。$e_1,m_1$ 分别
  是 2d 中的守恒动量与弦绕数，$\pi_1(\MH(G_{\mathrm{ad}}))=H^1(C,\ZZ)$ \pA\anchor{src p.115}。
\item \textbf{结论}：S-对偶在 2d 侧不是简单的"换 target"，而是同时交换：$\MH$ 的连通分支标签 $m_0$
  $\leftrightarrow$ B-field twist $b_{e_0}$，且交换 $G_{\mathrm{ad}}$ 与 $\Ldual_{\mathrm{ad}}$。
  因此 eigenbrane 必须理解成"所有 $e_0$ 的 twisted 对象之和"，单取一个 $e_0$ 是不闭的
  \pA\anchor{src p.116}。
\item \textbf{出口}：\S8.2 的"zerobrane 是电特征膜"只有在允许 $e_0$ 变化的 twisted 意义下才是干净的；
  \S11 的 $D$-模也必须是 twisted $D$-模。
\item \textbf{锚点索引}：p.110 (7.18)--(7.22) 分解；p.111--p.112 $e_0$ 作为平坦 B-field 的论证；
  p.113 膜支持分支 $\leftrightarrow$ B-field；p.114 Hitchin 纤维的 twisted duality；p.115 $e_1,m_1$
  与动量/绕数；p.116 两种量子化方式的交换。
\end{enumerate}
\textbf{可信度}：Künneth 分解与 $\xi(E)=(\Phi\times1)^*\zeta$ 是 A；$e_0$ 作为 B-field 的解释是 B
（作者在 p.112 给出逐项验证，但前提是低能 sigma 模型描述本身成立）\pB\anchor{src p.112}；$e_1,m_1$
的动量/绕数解释与最后"两种量子化的选择"是 C——原文用 "we can simply specify ... as we please" 给出
这种读法，属于作者的判断 \pC\anchor{src p.116}。
\end{kcard}

\begin{figure}[htbp]
\centering
\begin{tikzpicture}[font=\small,
  box/.style={draw=kwblue!65, fill=kwblue!4, rounded corners=2pt, align=center,
    inner sep=3pt, text width=4.0cm},
  lab/.style={font=\scriptsize, align=center}]
\node[box] (e0) {$e_0$：离散电 flux\\（$\in H^2(C)=\ZZ$）};
\node[box, right=1.0cm of e0] (b) {2d：$\MH$ 上的\\平坦 B-field $b_{e_0}$};
\node[box, right=1.0cm of b] (m0) {$m_0$：$\MH$ 连通分支\\（$\in H^2(C)=\ZZ$）};
\draw[{Latex[length=1.5mm]}-{Latex[length=1.5mm]}] (e0) -- node[above,lab]{$S$} (b);
\draw[{Latex[length=1.5mm]}-{Latex[length=1.5mm]}] (b) -- node[above,lab]{$S$} (m0);
\node[box, below=8mm of e0] (e1) {$e_1$：$C$ 上离散电 flux};
\node[box, right=1.0cm of e1] (p) {2d：守恒动量\\（$E_C=H^1(C,\ZZ)$）};
\node[box, right=1.0cm of p] (m1) {$m_1$：弦绕数};
\draw[{Latex[length=1.5mm]}-{Latex[length=1.5mm]}] (e1) -- node[above,lab]{$S$} (p);
\draw[{Latex[length=1.5mm]}-{Latex[length=1.5mm]}] (p) -- node[above,lab]{$S$} (m1);
\end{tikzpicture}
\caption{\S7.2 的翻译表：上排是 $H^2$ 层的 $(e_0,m_0)$（分别对应 B-field 与连通分支），下排是
$H^1$ 层的 $(e_1,m_1)$（分别对应守恒动量与绕数）。S-对偶在每一排内部做交换，同时交换 $G$ 与 $\Ldual$
\pA\anchor{src p.110}\anchor{src p.115}。}
\label{fig:geom-flux}
\end{figure}

% ======================================================================
\subsection{\S8 电特征膜：把 Wilson 的作用算出来（pp.116--123）}

\S8 是 S6 的收尾：\S8.1 把 \S6.4 的抽象作用 $XB=B\otimes V$ 算成具体公式 $U\mapsto U\otimes E_p(R)$，
并追踪它造成的 B-field 变化；\S8.2 用它证明 zerobrane 是电特征膜，从而把"电特征膜"具体化为
"$\pi_1(C)\to\LdualC$ 的表示"。本节是全篇第一次真正接触 GL 的主论断，但只走通了一个方向：电侧。
磁侧的对应物（type F 膜）要等到 \S9--\S11 才被证明有 Hecke 本征性质。

\S8 的可信度比较干净：核心公式 (8.7)--(8.9) 是 A，推导链条完整；需要打折扣的是"积分掉 $A_\Sigma$"
这一步（作者说对大的 $\mathrm{Im}\,\tau$ 在经典层面足够）以及省略的费米子项（作者说由拓扑对称唯一
决定）。\S8.2 末尾从 zerobrane 到 Hecke eigensheaf 的过渡是 B/C 混合：作者自己只说 "at least a
glimmering"。

\begin{kcard}{\S8.1 Wilson 算子如何作用于膜（pp.116--120）}
\begin{enumerate}[label=\textbf{\arabic*}]
\item \textbf{这一节的任务}：把 Wilson 线算子的 4d 表达式与 2d sigma 模型边界项的 Chan-Paton
  holonomy 逐项比较，求出它对膜的作用；再追踪这个作用如何改变离散电 flux（B-field）。
\item \textbf{输入}：\S6.1 的 Wilson 算子；\S6.4 的膜作用；\S7 的 universal bundle 与 twisted twist；
  \S4--\S5 的 $M=\Sigma\times C$ 紧化与 sigma 模型。
\item \textbf{关键构造}：取 $S=\gamma\times p$。把 $A=A+\ii\varphi$ 写成 pullback
  $\Phi_p^*(\hat A)$，其中 $\hat A$ 是 universal Higgs bundle 上由 (8.3)--(8.4) 决定的联络；于是
  Wilson 因子 (8.5) 与边界 Chan-Paton 因子 (8.6) 形式相同，得到
\begin{equation}
U\ \longmapsto\ U\otimes E_p(R)
\label{eq:geom-wilson-brane}
\end{equation}
  \pA\anchor{src p.118}\anchor{src p.119}。若 $E_p(R)$ 在表示 $R$ 的中心特征下 twisted，则 Wilson
  算子把背景 B-field 平移：$b\mapsto b+\theta_R(\zeta)$，等价地
\begin{equation}
e_0\ \longmapsto\ e_0+\theta_R ,
\label{eq:geom-e0-shift}
\end{equation}
  S-对偶下对应 't Hooft 算子改变拓扑类：$m_0\mapsto m_0+\xi(\LR)$ \pA\anchor{src p.119}\anchor{src p.120}。
\item \textbf{结论}：Wilson 算子的作用就是"张量上 $E_p(R)$"；它把膜移动一个 B-field 层，因此也把
  $e_0$ 平移。$E_p(R)$ 在每个复结构上全纯，所以这个作用保 (B,B,B) 型超对称——这正是 \S6.4 预言
  的 (B,B,B) 型线算子的几何实现 \pA\anchor{src p.121}。
\item \textbf{出口}：\S8.2 只需取 $B=B_x$（点支撑的 zerobrane），\eqref{eq:geom-wilson-brane} 立刻给出
  eigenbrane。
\item \textbf{锚点索引}：p.116--p.117 Wilson 算子与 Chan-Paton 边界项；p.117--p.118 用 universal
  bundle 表示 $\mathcal A$；p.119 (8.7) 张量积公式；p.119--p.120 B-field 平移与 $S$-对偶；p.121
  全纯性与 (B,B,B) 型。
\end{enumerate}
\textbf{可信度}：核心公式 (8.7)--(8.9) 是 A（推导完整，且与 \S6.4 的超对称预言一致）。打折处：
$\mathcal A_\Sigma$ 被经典地积分掉（作者说对大的 $\mathrm{Im}\,\tau$ 足够）\pB\anchor{src p.118}；
(8.3)--(8.6) 中省略的费米子项没有显式写出，作者断言它们由拓扑对称唯一决定，这是自然性论证
\pB\anchor{src p.119}。
\end{kcard}

\begin{kcard}{\S8.2 zerobrane $=$ 电特征膜（pp.121--123）}
\begin{enumerate}[label=\textbf{\arabic*}]
\item \textbf{这一节的任务}：证明点支撑的 zerobrane 是电特征膜，从而把 electric eigenbrane 具体化为
  "$\pi_1(C)\to\LdualC$ 的表示"，并说明 S-对偶把它送到 magnetic 侧的 type F 膜。
\item \textbf{输入}：\S8.1 的 \eqref{eq:geom-wilson-brane}；\S4.2（Hitchin 方程 $\leftrightarrow$ 平坦
  $\LdualC$-丛，复结构 $J$ 被拓扑超对称选出）；\S5.4--\S5.5（zerobrane 的 S-对偶是 Hitchin 纤维上的
  秩一平坦丛）。
\item \textbf{关键构造}：取 $x\in\MH$ 为光滑点，$B_x$ 是 $x$ 处的 zerobrane（Chan-Paton 为摩天层
  $U_x$）。任意向量丛在一点处平凡，故
\begin{equation}
U_x\otimes E_q(R)=U_x\otimes E_q(R)|_x ,
\label{eq:geom-zerobrane}
\end{equation}
  于是 $B_x$ 是电特征膜，特征值是向量空间 $E_q(R)|_x$ \pA\anchor{src p.121}。$q$ 变化时
  $E_q(R)|_x$ 是 $C$ 上平坦丛（复结构 $J$ 下 Hitchin 解就是平坦 $\LdualC$-丛）的纤维，与 \S6.4
  的预言一致 \pA\anchor{src p.121}\anchor{src p.122}。精确处理 B-field：电特征膜是"所有 $e_0$ 的
  twisted zerobrane 之和"；对偶地，magnetic eigenbrane 是"不同 $m_0$ 分支上的膜的之和"
  \pA\anchor{src p.122}。
\item \textbf{结论}：对每个同态 $\vartheta:\pi_1(C)\to\LdualC$（不可约、半稳定），点 $x(\vartheta)\in\MH(\Ldual)$
  处的 zerobrane 是电特征膜；S-对偶把它送到 $\MH(G)$ 中 Hitchin 纤维 $F$ 上的 type F 膜（秩一平坦
  Chan-Paton）。这是 GL 主论断的物理版\textbf{在一个方向上的雏形} \pA\anchor{src p.122}。
\item \textbf{出口}：\S9--\S11 要证明 type F 膜确实对应 Hecke 特征层：\S9--\S10 负责"磁算子 $=$
  Hecke"，\S11 负责"膜 $\to$ $D$-模"。没有后者，type F 膜只是几何对象，还不是 GL 的 Hecke eigensheaf。
\item \textbf{锚点索引}：p.121 (8.10) 与"任意丛在一点平凡"；p.121--p.122 平坦丛与 $J$ 复结构；
  p.122 B-field 与 twisted zerobrane 之和；p.122--p.123 从 zerobrane 到 GL 主论断的过渡（"a glimmering"）。
\end{enumerate}
\textbf{可信度}：计算 (8.10) 是 A。"不可约 $\vartheta$ 对应光滑点"是 B（依赖 \S4.2 的半稳定性与
不可约性条件）\pB\anchor{src p.122}。从"type F 膜"到"Hecke 特征层"的断言是 C：作者原文只说
"we can get at least a glimmering"，并把证明推给 \S9--\S11 \pC\anchor{src p.123}。
\end{kcard}

\begin{figure}[htbp]
\centering
\begin{tikzpicture}[font=\small,
  box/.style={draw=kwblue!65, fill=kwblue!4, rounded corners=2pt, align=center,
    inner sep=4pt, text width=4.3cm}]
\node[box] (w) {Wilson $W(R)$，$\Psi=\infty$\\ 作用：$U\mapsto U\otimes E_p(R)$};
\node[box, right=1.8cm of w] (t) {'t Hooft $T(\LR)$，$\Psi=0$\\ 作用：$m_0\mapsto m_0+\xi(\LR)$};
\node[below=5mm of w, font=\scriptsize, text width=4.3cm, align=center] (s1) {改变离散电 flux：$e_0\to e_0+\theta(R)$};
\node[below=5mm of t, font=\scriptsize, text width=4.3cm, align=center] (s2) {改变 $G$-丛拓扑类 $m_0$};
\draw[{Latex[length=1.5mm]}-{Latex[length=1.5mm]}] (w) -- node[above,font=\scriptsize]{$S$-对偶} (t);
\draw[-{Latex[length=1.4mm]}, kwgray] (w) -- (s1);
\draw[-{Latex[length=1.4mm]}, kwgray] (t) -- (s2);
\draw[-{Latex[length=1.4mm]}, kwred!60] (s1) to[out=-20,in=200] node[below,font=\scriptsize]{$S$-对偶交换两者} (s2);
\end{tikzpicture}
\caption{\S8.1 的 S-对偶方块：Wilson 在 $\Psi=\infty$ 平移 B-field（即 $e_0$），'t Hooft 在 $\Psi=0$
平移拓扑类 $m_0$；两者都是 (\S6.4) "算子改变膜"的具体化 \pA\anchor{src p.119}\anchor{src p.120}。}
\label{fig:geom-sdual-square}
\end{figure}

% ======================================================================
\subsection{\S9 't Hooft 算子 $=$ Hecke 算子（pp.123--148）}

\S9 开始兑现 \S6.4 的猜想：'t Hooft 算子的作用就是几何 Langlands 的 Hecke 算子。做法是把
$\Psi=0$ 的 4d 理论约化到 $W=\RR\times C$ 上的三维 Bogomolny 方程，再证明沿 $\RR$ 方向的演化只改变
全纯丛的 Hecke 修改型。本节输出 S7 的前一半——算子层面的等式；后一半（空间与算子代数的几何）留给
\S10。特别地，\S9.1 的 $U(N)$ 计算给出 $L_i\mapsto L_i\otimes\mathcal O(p_0)^{m_i}$ 这一行公式，它是
"物理算子 $=$ 数学修改"的全部内容。

\S9 的另一半工作是界定这一步的适用范围：约化到 2d 后是 $\MH$ 上复结构 $K$ 的 A-模型，BPS 构型是
全纯映射；由于 $\Sigma=\RR\times I$ 时解时间无关，"经典近似"就是常数映射的上同调（Floer 式）。作者
在 p.125--p.128 花了几页证明相关解确实时间无关（弱消失定理），这个证明只对 $(B,A,A)$ 型膜有效——
这是后面把 A-模型与膜型绑在一起的关键前提 \pB\anchor{src p.126}。

\begin{kcard}{\S9.1 静态 't Hooft 算子与 Hecke 修改（pp.128--134）}
\begin{enumerate}[label=\textbf{\arabic*}]
\item \textbf{这一节的任务}：把 4d 的 't Hooft 奇点在三维 Bogomolny 方程中逐分量解出来，证明它的
  效果是对全纯丛做 Hecke 修改；这是"物理算子 $=$ 数学算子"的第一个严格等式。
\item \textbf{输入}：\S6.2 的 't Hooft 算子 $T(\Lw)$；\S6.4 的静态线算子；S3（$M=\RR\times W$，
  $W=\RR^3=\RR\times C$）；$\Psi=0$ 的 A-模型。
\item \textbf{关键构造}：$z=x^2+\ii x^3$，静态算子位于 $y=y_0$、$z=p_0$。Bogomolny 方程
  $F=\star D\varphi_0$ 展开为
\begin{equation}
F_{z\bar z}=\frac{\ii}{2}D_y\varphi_0,\qquad
F_{yz}=\ii D_z\varphi_0,\qquad
F_{y\bar z}=-\ii D_{\bar z}\varphi_0 ,
\label{eq:geom-bogomolny-expand}
\end{equation}
  取 $A_y=0$ 时最后一式与规范变换相容：
  $\partial_yA_{\bar z}=-\ii D_{\bar z}\varphi_0$，即 $\partial_yD=-\ii[D,\varphi_0]$，故
  $E_y=E|_{\{y\}\times C}$ 的\textbf{全纯型在两次跨过 't Hooft 算子之间不变} \pA\anchor{src p.130}。
  $U(1)$：$L_+\otimes L_-^{-1}$ 在 $C\setminus\{p_0\}$ 上平凡，由面积分得
\begin{equation}
L_+\simeq L_-\otimes\mathcal O(p_0)^m ,
\label{eq:geom-hecke-u1}
\end{equation}
  即 $T(m;p_0)$ 的作用是 twist $\mathcal O(p_0)^m$ \pA\anchor{src p.131}\anchor{src p.132}。$U(N)$：
  取 $\Lw=(m_1\ge\cdots\ge m_N)$，近 $p_0$ 化到最大环使 $E_y=\oplus_iL_i$，则
\begin{equation}
L_i\ \longmapsto\ L_i\otimes\mathcal O(p_0)^{m_i},
\label{eq:geom-hecke-un}
\end{equation}
  与数学的 Hecke 算子逐项一致 \pA\anchor{src p.133}。$c_1(E)\mapsto c_1(E)+\sum_im_i$，拓扑不变当且
  仅当 $\sum_im_i=0$；$T(1),\dots,T(N-1)$ 与 $T(N)^{-1}$ 是自由生成元 \pA\anchor{src p.134}。
\item \textbf{结论}：'t Hooft 算子 $=$ Hecke 修改，且 $U(N)$ 的 Hecke 算子代理由基本 Hecke 操作生成；
  对 $PSU(N)$，$\Lw$ 取 $1/N$ 的分数值，$\mathcal O(p_0)^{1/N}$ 是 twisted 线丛，但差
  $m_i-m_j$ 为整数，故作用仍保持"所有 $L_i$ 的 twist 相同"这一事实 \pA\anchor{src p.135}。
\item \textbf{出口}：\S9.2 研究\textbf{所有可能的} Hecke 修改构成的空间 $Y(\Lw)$；\S10 把同一
  空间实现为 Bogomolny 模空间，从而能算 OPE。
\item \textbf{锚点索引}：p.128 Bogomolny 方程展开 (9.7)--(9.8)；p.130 全纯型不变性；(9.9)；p.131--p.132
  $U(1)$ 的 $L_+\simeq L_-\otimes\mathcal O(p_0)^m$；p.133 (9.13) 基本 Hecke 操作；p.134 维数公式
  出处；p.135 $PSU(N)$ 的分数权。
\end{enumerate}
\textbf{可信度}：推导与公式是 A（可由 \eqref{eq:geom-bogomolny-expand} 逐项核验）。"弱消失定理"是 B：
它只对 $(B,A,A)$ 型膜成立，作者在 p.126 明确说明这是本节结论的适用条件 \pB\anchor{src p.126}。
"算子代数 $\cong\mathrm{Rep}(\Ldual)$"在 p.134 是 C（原文说 "we expect"，依据是 S-对偶 + 4d TQFT）
\pC\anchor{src p.134}。
\end{kcard}

\begin{kcard}{\S9.2 Hecke 修改的空间 $Y(\Lw)$（pp.135--144）}
\begin{enumerate}[label=\textbf{\arabic*}]
\item \textbf{这一节的任务}：直接描述所有 Hecke 修改构成的空间 $Y(\Lw,p_0)$——用局部全纯截面
  给出模型、算出几个例子与维数、并构造自然紧化。这个空间将在 \S10 被实现为 Bogomolny 模空间。
\item \textbf{输入}：\S9.1 的局部公式；$C=\PP^1$、$p_0=0$、$G=U(N)$；$[122]$ 的 Schubert 胞腔理论。
\item \textbf{关键构造}：设 $s_1,\dots,s_N$ 是 $E_-$ 在 $p_0$ 附近的局部全纯截面、在 $0$ 处线性无关，
  则 $E_-=\oplus_iL_i$（$L_i$ 由 $s_i$ 生成）；$T(\Lw)$ 把 $E_+$ 的截面写成
  $\sum_i g_i(z)s_i(z)z^{-m_i}$。例子：$\Lw=(1,0)$ 时 $Y=\PP^1$；$\Lw=(2,0)$ 时参数
  $(u,v,b=vu'-uv')$ 给加权射影空间 $\PP(1,1,2)$，在 $(0,0,1)$ 处有 $A_1$ 型奇点 $C^2/\ZZ_2$，
  其紧化必须加入类型 $(1,1)$，即"张量 $\mathcal O(p_0)$" \pA\anchor{src p.136}\anchor{src p.138}。
  解消：把 $\PP(1,1,2)$ 写成 $\CC^4$ 的超凯勒商，矩映射
  $|u|^2+|v|^2+2|b|^2=1$、$|b|^2+|b'|^2=d$；$d<1/2$ 给 Hirzebruch 面 $Y=\PP(\mathcal O(2)\oplus\mathcal O)$，
  $d=1/2$ 退化回 $\PP(1,1,2)$；形变 $a_1a_3-a_2^2=\epsilon b^2$ 给 $\PP^1\times\PP^1$
  \pA\anchor{src p.140}\anchor{src p.141}。一般 $\Lw$：若 $m_i\ge m_j+2$，紧化必须加入
  $(m_i,m_j)\to(m_i-1,m_j+1)$（"associated weights"）\pA\anchor{src p.143}；minuscule 表示
  （$\Lw(k)$）对应的 $Y=\mathrm{Gr}(k,N)$ 光滑紧 \pA\anchor{src p.144}；维数
\begin{equation}
\Delta_{\Lw}=\sum_{i<j}(a_i-a_j).
\label{eq:geom-hecke-dim}
\end{equation}
\item \textbf{结论}：$Y(\Lw)$ 是有限维复代数簇，一般奇异且非紧；其自然紧化 = 加入所有
  associated weights 的修改；minuscule 表示是"光滑紧"的例外。
\item \textbf{出口}：\S9.3 用仿射 Grassmannian 解释"紧化为什么必须加入这些类型"；\S10.2 用单极子成泡
  给出同一件事的物理解释；\S10.4 用 $Y$ 的奇点计算 OPE 的系数空间。
\item \textbf{锚点索引}：p.136--p.137 局部截面生成与 (9.16)--(9.17)；p.138--p.139 加权射影空间与紧化；
  p.140--p.141 超凯勒商、Hirzebruch 面、形变；p.143 associated weights；p.144 minuscule 与
  Grassmannian；p.144 维数公式 (9.25)。
\end{enumerate}
\textbf{可信度}：局部模型、例子与维数公式是 A（原文计算完整，可自行核验 (9.25)）。"自然紧化包含
associated weights"是 A$^-$：作者在 \S9.2 用 $U(2)$ 族与退化论证给出，完整的证明在 \S9.3（轨道闭包）
与 \S10.2（单极子成泡）\pA\anchor{src p.143}。单极子成泡的物理图像本身是 B \pB\anchor{src p.155}。
\end{kcard}

\begin{kcard}{\S9.3 仿射 Grassmannian 与 Schubert 胞腔（pp.145--148）}
\begin{enumerate}[label=\textbf{\arabic*}]
\item \textbf{这一节的任务}：把 $Y(\Lw)$ 放进仿射 Grassmannian 的框架，证明它就是一个 Schubert
  胞腔，从而得到自然紧化（Schubert cycle）。作者声明本文不依赖这个框架，但它解释了紧化的来历。
\item \textbf{输入}：\S9.2；$G=U(N)$；$[122]$ 第 8 章的标准构造。
\item \textbf{关键构造}：Hecke 修改是带同构 $\sigma:E\simeq E_-$（$p_0$ 外）的丛，等价类空间
  $=\mathrm{GL}(N,\mathcal X)/\mathrm{GL}(N,\mathcal O)=\mathrm{Gr}_N$，其中 $\mathcal X$ 是 $C^*$ 上
  可亚纯延拓的全纯函数环；也可用形式版本 $\mathrm{GL}(N,\CC((z)))/\mathrm{GL}(N,\CC[[z]])$
  \pA\anchor{src p.145}\anchor{src p.146}。权 $\Lw$ 的 Hecke 修改 $=$ 点
  $g_{ji}(z)=z^{-m_j}\delta_{ji}$ 在 $\mathrm{GL}(N,\mathcal O)$ 下的轨道；所有轨道都这样得到
  $[122]$，故 $\mathrm{Gr}_N$ 按 Schubert 胞腔分层 \pA\anchor{src p.147}。$\mathrm{Gr}_N(i)$
  （在 $z=0$ 处极阶 $\le i$ 的子集）有限维紧且 $\mathrm{GL}(N,\mathcal O)$-不变，故每个轨道有限维；
  紧化 $=$ 轨道闭包 $=$ Schubert cycle，由低维轨道组成 \pA\anchor{src p.148}。一般 $G$：
  $\mathrm{Gr}_G=G_\CC(\mathcal X)/G_\CC(\mathcal O)$，$Y(\rho)$ 是 $\rho_\CC:\CC^*\to G_\CC$
  所定点的轨道 \pA\anchor{src p.149}。
\item \textbf{结论}：$Y(\Lw)$ $=$ 仿射 Grassmannian 中的 Schubert 胞腔；其闭包（Schubert
  cycle）给出自然紧化。这从纯几何上解释了 \S9.2 中"必须加入 associated weights"的现象。
\item \textbf{出口}：\S10.2 用单极子成泡给同一现象以物理机制；\S10.4 用 Schubert 簇的交截上同调把
  OPE 系数精确化（Lusztig--Ginzburg--Mirkovi\'c--Vilonen）。
\item \textbf{锚点索引}：p.145--p.146 $\mathrm{GL}(N,\mathcal X)/\mathrm{GL}(N,\mathcal O)$ 与形式版本；
  p.147 (9.29)--(9.31) 轨道描述；p.148 $\mathrm{Gr}_N(i)$ 有限维、Schubert cycle；p.149 一般 $G$ 与
  $\mathrm{Gr}_G$。
\end{enumerate}
\textbf{可信度}：全部构造是 A（标准数学，原文明确引用 $[122]$）。作者在 p.145 明说"本文将不真正
使用这个框架"，因此这里没有物理断言可打折；唯一带判断色彩的是把"Schubert cycle $=$ 紧化"与物理的
"混合"对应起来，属于 B \pB\anchor{src p.149}。
\end{kcard}

\begin{figure}[htbp]
\centering
\begin{tikzpicture}[font=\small, >={Latex[length=1.8mm]}]
% ---------- (a) 局部 Hecke 修改 ----------
\node[font=\small] at (-4.0,3.05) {\textbf{(a)} 局部：$\Lw=(m_1,\dots,m_N)$};
\node[align=center, font=\scriptsize] (em) at (-6.0,1.95) {$E_-=\bigoplus_iL_i$};
\node[align=center, font=\scriptsize] (ep) at (-6.0,1.05) {$E_+=\bigoplus_iL_i\otimes\mathcal O(p_0)^{m_i}$};
\draw[->] (em.south) -- (ep.north) node[midway,right,font=\scriptsize]{$T(\Lw)$};
\draw[thick, kwblue!70] (-7.9,2.15) .. controls (-7.3,2.42) .. (-6.7,2.15);
\fill (-7.3,2.28) circle (1.6pt);
\node[font=\scriptsize, below] at (-7.3,2.08) {$p_0$};
\node[align=left, font=\scriptsize, text width=3.6cm] at (-2.9,2.35)
  {截面 $s_i\mapsto s_i/z^{m_i}$\\[2pt]（$E_+$ 只依赖 $s(0)$ 的射影类）};
% ---------- (b) U(1) 谱流 ----------
\node[font=\small] at (3.4,3.05) {\textbf{(b)} $U(1)$：谱流};
\draw[thick] (1.7,1.0)--(1.7,2.1);
\draw[thick] (4.7,1.0)--(4.7,2.1);
\node[font=\scriptsize, below] at (1.7,0.95) {$L_-$};
\node[font=\scriptsize, below] at (4.7,0.95) {$L_+$};
\draw[->] (1.7,1.98) to[bend left=20] node[above,font=\scriptsize]{$\otimes\,\mathcal O(p_0)^m$} (4.7,1.98);
\node[font=\scriptsize] at (3.2,0.15) {$q_+-q_-=m$（图 14 的 cobordism）};
\end{tikzpicture}
\caption{Hecke 修改的两种图像：局部是"每个分量乘一个点幂"（\eqref{eq:geom-hecke-un}），整体是
$U(1)$ 情形的谱流 $q_+-q_-=m$。一般 $G$ 的 Hecke 修改空间就是这个局部模型的所有可能实现的模空间
\pA\anchor{src p.133}\anchor{src p.132}。}
\label{fig:geom-hecke}
\end{figure}

% ======================================================================
\subsection{\S10 Bogomolny 方程与 Hecke 修改空间的几何（pp.148--183）}

\S9 证明了"算子 $=$ 修改"，但没有解释修改空间为什么长成那样、奇点从哪里来、OPE 系数怎么算。
\S10 的任务是把 $Y(\Lw)$ 实现为 Bogomolny 方程带奇点的解模空间 $Z$，然后借助微分几何与指标
定理把 \S9 的代数观察逐条翻译成几何事实。五个小节的逻辑是：先定边界条件（\S10.1），用指数定理算
维数；再用单极子成泡解释紧化与奇点（\S10.2）；用 Kähler 结构/$ $矩映射把 $Z$ 等同于 Hecke 修改空间
并解释边界条件的自然性（\S10.3）；最后用 $Z$ 的奇点与 Schubert 簇的上同调算出 OPE（\S10.4）；
\S10.5 放开时间反演，给出 extended Bogomolny 方程，这是让 't Hooft 算子作用于一般 A-膜的必需工具。

\begin{kcard}{\S10.1 边界条件与指数定理（pp.148--154）}
\begin{enumerate}[label=\textbf{\arabic*}]
\item \textbf{这一节的任务}：在 $W=I\times C$ 上给 Bogomolny 方程配一组椭圆边界条件，使解模空间
  $Z$ 恰好是 Hecke 修改空间；再用指数定理算出 $\dim Z$，并与 \S9.2 的 \eqref{eq:geom-hecke-dim} 对照。
\item \textbf{输入}：\S9 的 't Hooft 奇点；\S3--\S4 的场论设置；指标理论 $[107,108]$、Kronheimer $[110]$、
  Pauly $[124]$。
\item \textbf{关键构造}：左端 $C_-$：固定 $A_-$（Dirichlet），$\varphi_0$ 自由；右端 $C_+$：固定
  $\varphi_0=0$，$A$ 满足协变 Neumann $\iota_nF=0$；规范变换在左端平凡、右端不平凡。这样 $E_-$ 是
  输入、$E_+$ 是输出（图 15(a)）。Hecke correspondence 的版本是两端都固定 $A$（图 15(b)），对应
  非椭圆情形。线性化给出椭圆复形
\begin{equation}
0\to\Omega^0(\mathrm{ad}E)\xrightarrow{d_1}\Omega^1(\mathrm{ad}E)\oplus\Omega^0(\mathrm{ad}E)
\xrightarrow{d_2}\Omega^2(\mathrm{ad}E)\to0,
\label{eq:geom-elliptic}
\end{equation}
  $\text{index}=\dim Z$（在光滑点）。无 't Hooft 算子时 $Z$ 是零维点：由
  $\int_W\Tr(F-\star D\varphi_0)\wedge\star(F-\star D\varphi_0)=0$ 与边界项消失得
  $F=D\varphi_0=0$ \pA\anchor{src p.153}\anchor{src p.154}。每个 $T(\rho)$ 对维数的贡献只依赖
  $\Ldual$ 的表示，Pauly/Kronheimer 给出
\begin{equation}
\Delta_\rho=\sum_kc_k,
\label{eq:geom-pauly}
\end{equation}
  其中 $\rho:U(1)\to G$ 在伴随表示下分解为权 $c_k,-c_k$；单算子时 $Z=Y(\rho)$，\eqref{eq:geom-pauly}
  与 \eqref{eq:geom-hecke-dim} 一致 \pA\anchor{src p.155}\anchor{src p.156}。
\item \textbf{结论}：$Z$ 的维数 $=\sum_i\Delta_{\rho_i}$，与 Hecke 修改空间的维数相符。边界条件是
  "使 Bogomolny 方程成为全纯方程 $+$ 矩映射消失"的最简选择（\S10.3 给正式解释）。
\item \textbf{出口}：\S10.2 用成泡解释 $Z$ 的奇点与紧化；\S10.3 用它证明 $Z$ 的复结构就是 Hecke
  修改空间。
\item \textbf{锚点索引}：p.149--p.150 图 15 与两类问题；p.151 椭圆边界条件；p.152 指数复形 (10.1)；
  p.153--p.154 无算子时 $Z$ 是点；p.154--p.156 Pauly 的 $\Delta_\rho$ 与 (10.4)--(10.5)。
\end{enumerate}
\textbf{可信度}：边界条件与指数定理是 A（引 $[124]$ 的定理并复述证明思路）。"$H^2$ 在光滑点消失"
是 A（按指标理论）；"整体上 $Z$ 是 $\dim Z$ 维流形"是 B（作者只断言在光滑部分，奇异部分留给
\S10.2）。"对一般 $G$ 的 $\Delta_\rho$ 公式"是 A$^-$：原文说是 Pauly 证明的直接推广，但没有给出
完整细节 \pB\anchor{src p.156}。
\end{kcard}

\begin{kcard}{\S10.2 单极子成泡（pp.155--158）}
\begin{enumerate}[label=\textbf{\arabic*}]
\item \textbf{这一节的任务}：解释 \S9.2 中"紧化必须加入 associated weights"的物理机制——单极子成泡
  (monopole bubbling)，并算出成泡如何改变 't Hooft 算子的权。
\item \textbf{输入}：\S9.2 的紧化；Kronheimer 的 $F$-等变瞬子 $\leftrightarrow$ Bogomolny 解；ADHM 构造
  $[125]$ 与 $[126]$。
\item \textbf{关键构造}：$F=U(1)$ 以 $(z_1,z_2)\mapsto(e^{\ii\theta}z_1,e^{\ii\theta}z_2)$ 作用于
  $C^2$，$C^2/F\simeq\RR^3$，映射为 $\vec x=z\vec\sigma\bar z$；$F$-等变瞬子在原点的纤维
  $E_0$ 决定 't Hooft 算子的权 \pA\anchor{src p.155}。瞬子方程尺度不变 $\Rightarrow$ 瞬子可缩成点
  （"bubbling"），在三维图像中就是单极子成泡：一个量子化电荷集中到 't Hooft 算子所在点并消失
  \pA\anchor{src p.157}。用 ADHM 复形 (10.7) 与特征公式
\begin{equation}
\chi(E_0')=\chi(E_0)+\chi(H_1)\,(2-e^{\ii\theta}-e^{-\ii\theta})
\label{eq:geom-bubbling}
\end{equation}
  得到权的变化 (10.13)：$(m+1,m-1)\to(m,m)$；$H_1$ 的秩为 $s$ 且权为 $m,\dots,m+s-1$ 时
  $(m+s,m-1)\to(m+s-1,m)$ \pA\anchor{src p.158}\anchor{src p.159}。奇点的局部超凯勒结构来自
  $F$-等变瞬子模空间（继承 $\RR^4$ 的超凯勒结构）\pA\anchor{src p.159}。
\item \textbf{结论}：associated weights 的加入 $=$ 单极子成泡；$Y(\Lw)$ 的奇点是超凯勒奇点，
  可形变也可解消（例：$\PP(1,1,2)$ 的 $A_1$ 奇点）。这解释了 \S9.2 的紧化为何"恰好"只加入
  associated types。
\item \textbf{出口}：\S10.3 把"沿 $I$ 方向移动 $=$ 解消、沿 $C$ 方向移动 $=$ 形变"讲清楚，并用于
  OPE 极限；\S10.4 用奇点算系数。
\item \textbf{锚点索引}：p.155 $F$-等变瞬子与 Kronheimer 对应；p.156--p.157 成泡与电荷消失；
  p.157--p.158 ADHM 复形与 (10.12)--(10.13)；p.159 超凯勒结构来源。
\end{enumerate}
\textbf{可信度}：ADHM 计算与权变化公式是 A（可核验 (10.13)）。"单极子成泡"的图像是 B：它是
Kronheimer/Pauly 数学结论的物理读法，作者用 "we might call" 引入这个术语 \pB\anchor{src p.157}；
"对一般 $G$ 的推广可用嵌入的 $SU(2)$ 生成"是 B/C（作者说 "we will see that bubbling changes the
weight in precisely the expected fashion"，对例外群只能类比）\pC\anchor{src p.157}。
\end{kcard}

\begin{kcard}{\S10.3 Kähler 结构、矩映射与 OPE 极限（pp.159--165）}
\begin{enumerate}[label=\textbf{\arabic*}]
\item \textbf{这一节的任务}：给场空间配 Kähler 结构，证明 $Z$ 的复结构 $=$ Hecke 修改空间；由此
  解释边界条件的自然性，并算出两个算子靠近时的极限 $Z^*=Z(\Lw+\Lw')$。
\item \textbf{输入}：\S10.1 的边界条件；\S10.2 的成泡。
\item \textbf{关键构造}：取 $A=A_z\dd z+A_{\bar z}\dd\bar z+A_y\dd y$，Bogomolny 方程写成
\begin{equation}
[D_y-\ii\varphi_0,D_z]=0,\qquad F_{23}=hD_y\varphi_0 .
\label{eq:geom-bog-kahler}
\end{equation}
  在场空间上取复结构 $I$（坐标 $A_z$、$B=A_y-\ii\varphi_0$）与辛形式
  $\omega=\int_W\Tr(\delta A_2\wedge\delta A_3-h\,\delta A_1\wedge\delta\varphi_0)$，合成 Kähler
  结构 \pA\anchor{src p.160}\anchor{src p.161}。矩映射
\begin{equation}
\mu(\epsilon)=\int_W\Tr\big(\epsilon(F_{23}-hD_y\varphi_0)\big)+\int_{\partial W}\Tr(\epsilon\varphi_0)h ,
\label{eq:geom-moment}
\end{equation}
  $\mu=0$ 给出 \eqref{eq:geom-bog-kahler} 的第二式，并要求边界上 $\Tr\epsilon\varphi_0=0$；两种边界选择
  （$\epsilon=0$ 或 $\varphi_0=0$）正好给出图 15(a) 的左右边界 \pA\anchor{src p.162}。于是辛商
  $=$ 复商 $Z=C_0/G_\CC$（无稳定性问题，因一端规范平凡）；$C_0/G_\CC$ 的每一点由"跨算子做 Hecke
  修改"唯一决定，故 $Z$ 的复结构 $=Y$ \pA\anchor{src p.163}。
\item \textbf{结论}：$Z$ 有自然 Kähler 结构；沿 $I$ 方向滑动插入点只改变辛结构、不改变复结构；沿
  $C$ 方向移动只改变复结构、不改变辛结构（因为 $\omega$ 只用面积形式）。令 $p'\to p$ 时
  $Z(\Lw,p;\Lw',p')\to Z(\Lw+\Lw',p')$，展开为
\begin{equation}
T(\Lw,p)\cdot T(\Lw',p')\to T(\Lw+\Lw',p')
+\sum_i c_i\,T(\Lw_i,p'),
\label{eq:geom-ope-limit}
\end{equation}
  其中 $\Lw_i$ 是 $\Lw+\Lw'$ 的 associated weights；维数匹配 \eqref{eq:geom-pauly}
  \pA\anchor{src p.165}\anchor{src p.166}。
\item \textbf{出口}：\S10.4 用奇点的交截上同调把 \eqref{eq:geom-ope-limit} 的系数 $c_i$ 提升为向量空间
  $N_\alpha$；\S10.5 放开时间反演。
\item \textbf{锚点索引}：p.160 (10.14)--(10.17) 复结构与辛形式；p.161--p.162 矩映射与边界条件；
  p.162--p.163 辛商 $=$ 复商与 $Z$ 的复结构；p.163--p.164 $D_I$ 作用与形变/解消；p.165--p.166
  (10.20) OPE 极限与 associated weights。
\end{enumerate}
\textbf{可信度}：Kähler 结构的微分几何计算是 A。辛商 $=$ 复商的稳定性论证是 B（作者说"因为一端
规范平凡，故 $G_\CC$ 自由作用、轨道闭，无稳定性问题"，论证合理但只作断言）\pB\anchor{src p.162}。
"$c_i$ 作为关联权重的系数"是 A（由维数匹配与 $Z$ 的极限确定）；但"OPE 的完整内容"仍待 \S10.4，
作者在 p.167 明说这是难点 \pA\anchor{src p.166}。
\end{kcard}

\begin{kcard}{\S10.4 't Hooft 算子的 OPE（pp.167--178）}
\begin{enumerate}[label=\textbf{\arabic*}]
\item \textbf{这一节的任务}：确定 \eqref{eq:geom-ope-limit} 的系数——证明它们是 $\Ldual$ 表示张量积的重数
  空间 $N_\alpha$，从而把 \S6.4 的猜想 (6.29) 变成严格陈述。
\item \textbf{输入}：\S6.4 的 (6.29)；\S10.3 的模空间与极限；A-模型：态空间 $=H^*(Z)$；数学文献
  Lusztig $[96]$、Ginzburg $[97]$、Mirkovi\'c--Vilonen $[98]$。
\item \textbf{关键构造}：记
\begin{equation}
\LR\otimes\LR'=\bigoplus_{\alpha\in R}N_\alpha\otimes\LR_\alpha ,
\qquad n_\alpha=\dim_\CC N_\alpha .
\label{eq:geom-tensor-decomp}
\end{equation}
  静态 Wilson 算子的精确 OPE 是"向量空间系数"版本
  $W(\LR,p)W(\LR',p')\to\bigoplus_\alpha N_\alpha\otimes W(\LR_\alpha,p')$；
  S-对偶期望 't Hooft 算子同样 (10.23) \pA\anchor{src p.169}\anchor{src p.170}。若把"时间"取成圆，
  迹把 $N_\alpha$ 换成 $n_\alpha$，这就是圈算子 OPE 的整数系数。取 brane $B,B'$ 为 $(B,A,A)$ 型
  （$U$：固定全纯型、$\phi$ 任意；$U'$：$\phi=0$），则态空间
  \begin{equation}
H(\LR_1,\dots,\LR_k)=H^*\big(Z(\LR_1,\dots,\LR_k)\big),
\label{eq:geom-states}
\end{equation}
  且 $Z$ 拓扑上是乘积，故 $H$ 的乘法与 $\LR$ 的张量积同构；由 $[127]$，$H(\LR)\cong\LR$
  \pA\anchor{src p.172}\anchor{src p.173}\anchor{src p.174}。奇点的处理：令
  $\pi:Z(\Lw,z,y;\Lw',z,y')\to Z(\Lw+\Lw')$ 为 blowdown（$\pi$ 在光滑部分
  是同构），则
\begin{equation}
\pi_*\big(IC(\Lw)\star IC(\Lw')\big)
=IC(\Lw+\Lw')\oplus\bigoplus_\alpha\widetilde N_\alpha\otimes IC(\Lw_\alpha),
\label{eq:geom-ic}
\end{equation}
  首项系数是 1 维（对应 (6.29) 的"最高权相加为主"），$\widetilde N_\alpha$ 是其余系数空间；
  $[96]$--$[98]$ 的定理证明 $\widetilde N_\alpha\cong N_\alpha$ \pA\anchor{src p.178}。
\item \textbf{结论}：'t Hooft 算子的 OPE 系数空间就是表示张量积的重数空间 $N_\alpha$。因此 \S6.4 的
  猜想在数学上是真的；但要注意：物理侧只推出 \eqref{eq:geom-ope-limit} 的\textbf{首项}（主导权），
  把 $\widetilde N_\alpha$ 等同到 $N_\alpha$ 是数学定理，不是本文物理论证的结论。
\item \textbf{出口}：\S11 用 D-模把这套几何再翻译一次；\S12 把 $B,B'$ 的选择整理成检查清单。
\item \textbf{锚点索引}：p.167--p.169 (10.21)--(10.23) 张量积与 OPE 形式；p.169--p.170 态空间
  $=\;H^*(Z)$；p.172--p.174 乘积结构与 $H(\LR)\cong\LR$；p.178 (10.33)--(10.34) 交截上同调
  与 $[96]$--$[98]$。
\end{enumerate}
\textbf{可信度}：OPE 的形式与态空间恒等式是 A（物理上清楚，代数上可核验）。"$N_\alpha$ 的确定"
是 A，但依据是数学定理 $[96]$--$[98]$，不是本文物理推导；作者在 p.169 明说"hard part of the proofs"
在数学文献 \pA\anchor{src p.169}。中间有一段 C：作者承认物理解释 "$H(\LR)\cong\LR$" 的
$\Ldual$ 作用"from a physical point of view is mysterious"，并猜测理解 $B,B'$ 的 S-对偶能澄清
\pC\anchor{src p.174}。
\end{kcard}

\begin{kcard}{\S10.5 Extended Bogomolny 方程（pp.179--183）}
\begin{enumerate}[label=\textbf{\arabic*}]
\item \textbf{这一节的任务}：去掉时间反演不变性，写出 extended Bogomolny 方程；这是让 't Hooft 算子
  作用于一般 A-膜（特别是 Hitchin 纤维与 \S11 的 c.c. 膜）的必要工具。
\item \textbf{输入}：\S9 开头的 BPS 方程；\S10.1--\S10.4；$[26]$（详细计算）。
\item \textbf{关键构造}：写 $\varphi=\varphi_0\dd x^0+\pi^*\tilde\varphi$、
  $A=A_0\dd x^0+\tilde A$，时间无关的 BPS 方程 (10.35)；用与 (9.4) 同样的正定论证得
  $\varphi_1=A_0=0$；取 $A_y=0$、度量 $\dd s^2=h|\dd z|^2+\dd y^2$，方程化为
\begin{equation}
D_z\varphi_z=0,\quad
\partial_yA_{\bar z}=-\ii D_{\bar z}\varphi_0,\quad
\partial_y\varphi_z=-\ii[\varphi_z,\varphi_0],\quad
F_{z\bar z}-[\varphi_z,\varphi_{\bar z}]=\frac{\ii h}{2}\partial_y\varphi_0 .
\label{eq:geom-extended}
\end{equation}
  第一式说 $\varphi_z\dd z$ 是 $\mathrm{ad}E\otimes K_C$ 的全纯截面，故 $(E,\varphi)$ 是 Higgs
  丛；后两式说 $y$ 演化只差规范变换，故 $\MH$ 中的点与 $y$ 无关 \pA\anchor{src p.181}。
\item \textbf{结论}：跨过 't Hooft 算子时，$(E,\varphi)$ 做 Hecke 修改；与丛的情形不同，全纯截面
  $\varphi$ 不能在一点被任意修改，故 Higgs 丛的 Hecke 修改一般是有限集，且保 Hitchin 纤维
  \pA\anchor{src p.182}\anchor{src p.183}。取 $G=U(N)$、$\Lw=(m_1,\dots,m_N)$、$E_+$ 由
  $t_i=(z-z_0)^{-m_i}s_i$ 生成，则正规性条件为
\begin{equation}
\varphi\cdot t_i=\sum_j(z-z_0)^{m_j-m_i}f_{ij}\,t_j ,
\label{eq:geom-higgs-hecke}
\end{equation}
  即 $\varphi$ 在 $Y(\Lw,z_0)$ 上定义向量场，零点 $=$ 可能的 Hecke 修改；$\varphi(z_0)$
  正则半单时这些零点就是特征方向 \pA\anchor{src p.182}\anchor{src p.183}。
\item \textbf{出口}：\S11 的 c.c.\ A-膜支撑在 $\MH$ 的余各向同性子流形上；extended Bogomolny 方程
  保证它不被 't Hooft 算子移出支撑，从而 magnetic eigenbrane 的论证才能进行。这也是 $[26]$ 中
  "不用 S-对偶直接证明"的路线。
\item \textbf{锚点索引}：p.180 (10.35) 与消失定理；p.181 (10.36) extended 方程与 Higgs 丛解释；
  p.182 (10.37)--(10.39) 全纯条件与 $f_{ij}$；p.182--p.183 零点、特征值、Hitchin 纤维不变性。
\end{enumerate}
\textbf{可信度}：方程 \eqref{eq:geom-extended} 与推导是 A。“$\varphi$ 消去多余参数 $\Rightarrow$ Hecke 修改
有限”是 B：作者用 "strongly suggests" 措辞，并只在 $\varphi(z_0)$ 正则半单的情形断言 \pB\anchor{src p.181}。
"支撑在 Hitchin 纤维上的膜是 magnetic eigenbrane 的一步"是 C：作者明确说完整证明在 $[26]$
\pC\anchor{src p.183}。
\end{kcard}

\begin{figure}[htbp]
\centering
\begin{tikzpicture}[font=\small,
  bnd/.style={font=\scriptsize, align=center, text width=2.9cm}]
% 上：W = I x C 条带
\draw[thick, fill=kwblue!4] (0,0) rectangle (6.2,2.6);
\node[bnd] at (-1.55,1.35) {$C_-$：固定 $A_-$\\（Dirichlet），$\varphi_0$ 自由};
\node[bnd] at (7.75,1.35) {$C_+$：$\varphi_0=0$，\\$\iota_nF=0$（Neumann）};
\foreach \x in {2.0,3.2,4.4} { \fill (\x,1.3) circle (2pt); }
\node[font=\scriptsize, above] at (2.0,1.45) {$T(\Lw_1)$};
\node[font=\scriptsize, above] at (3.2,1.45) {$T(\Lw_2)$};
\node[font=\scriptsize, above] at (4.4,1.45) {$T(\Lw_3)$};
\node[font=\scriptsize] at (3.1,0.45) {$W=I\times C$};
\draw[-{Latex[length=1.6mm]}] (-0.35,1.3) -- (0,1.3);
\draw[-{Latex[length=1.6mm]}] (6.2,1.3) -- (6.55,1.3);
\node[font=\scriptsize, below] at (0.9,0.05) {$E_-$ 输入};
\node[font=\scriptsize, below] at (5.3,0.05) {$E_+$ 输出};
% 下：两个算子靠近，退化为单算子
\node[font=\scriptsize, anchor=west] at (-1.55,-0.85)
  {当 $p'\to p$：两个算子退化成一个（单极子成泡）};
\fill (0.4,-1.75) circle (2pt);  \node[font=\scriptsize, above left]  at (0.35,-1.7) {$p$};
\fill (1.2,-1.75) circle (2pt);  \node[font=\scriptsize, above right] at (1.25,-1.7) {$p'$};
\draw[-{Latex[length=1.6mm]}] (1.55,-1.75) -- (2.2,-1.75);
\fill[kwred!70] (2.45,-1.75) circle (2pt);
\node[font=\scriptsize, anchor=west] at (2.75,-1.75) {$Z(\Lw+\Lw',p)$（单算子）};
\node[font=\scriptsize, anchor=north west] at (0.15,-2.1) {$Z(\Lw,p)\times Z(\Lw',p')$};
\end{tikzpicture}
\caption{\S10.1--\S10.3 的几何设置：条带 $W=I\times C$ 两端放 $(B,A,A)$ 型膜，中间插 't Hooft 算子；
解模空间 $Z$ 的复结构就是 Hecke 修改空间。两个算子靠近时 $Z$ 退化为单算子模空间，其奇异 locus
就是单极子成泡的位置，也是 OPE 系数所在 \pA\anchor{src p.150}\anchor{src p.166}。}
\label{fig:geom-bogomolny}
\end{figure}

% ======================================================================
\subsection{\S11 A-膜与 D-模（pp.184--206）}

\S11 补上 \S8.2 留下的最后一块缺口。几何 Langlands 的陈述是：$\pi_1(C)\to\LdualC$ 的同态对应
$\MC$ 上既是 Hecke 特征层又是 $D$-模的对象。左边我们已经认识——电特征膜（\S8），其对偶是
Hitchin 纤维上的 type F 膜；但右边在物理侧完全缺席：type F 膜是 $\MH$ 上的 $(A,B,A)$ 膜，
而不是 $\MC$ 上的 $D$-模。\S11 的断言是：\textbf{$\MH$ 上复结构 $K$ 中的每个 A-膜都自然给出
$\MC$ 上的一个 $D$-模}，机制是典范余各向同性膜（c.c.\ 膜）——它是 space-filling 的 $(A,B,A)$ 膜，
它的开弦代数恰是全纯微分算子层，而任何其它 A-膜都是这个代数的模 \pA\anchor{src p.184}。

读这一节要抓住三个转译：(i) c.c.\ 膜 $\leftrightarrow$ 微分算子层 $\Dmod_{K_{\MC}^{1/2}}$
（\S11.1）；(ii) 任意 A-膜 $\leftrightarrow$ 该层的模，即 $D$-模（\S11.2）；(iii) 把 $B$-场参数
$\Psi$ 开起来时，同一构造变成 twisted $D$-模，并在 $\Psi=h$ 时退化为普通 $D$-模（\S11.3）。
三个转译的严格程度递减：第一段是逐项可算的开弦代数，第二段依赖开弦代数只能"部分层化"
（只能对 $T^*V$ 型开集定义），第三段的全局粘接修正只是二阶摄动论证。这决定了 \S11 的证据分级。

\begin{kcard}{\S11.1 典范余各向同性膜与开弦代数（pp.185--196）}
\begin{enumerate}[label=\textbf{\arabic*}]
\item \textbf{这一节的任务}：构造支撑整个 $\MH$ 的余各向同性 A-膜 $B_{c.c.}$，并计算它的
  $(B,B)$ 开弦代数——结论是全纯微分算子层；这是 \S11.2 把 A-膜翻成 $D$-模的代数基础。
\item \textbf{输入}：\S4.1 的三个辛形式 $\omega_I,\omega_J,\omega_K$；\S5.3 的膜型表；$[80]$ 的
  余各向同性膜；$[41]$ 的 A-膜 $\leftrightarrow$ 非交换 B-膜。
\item \textbf{关键构造}：space-filling 的余各向同性 A-膜要求其 Chan-Paton 曲率满足
\begin{equation}
(\omega^{-1}F)^2=-1 ,
\label{eq:geom-cc-condition}
\end{equation}
  此时 $\widehat N=\omega^{-1}F$ 是可积复结构，而 $\omega-\ii F$ 是它上面的全纯辛形式；故这类膜只在
  $\omega$ 为全纯辛形式实部时存在 \pA\anchor{src p.185}。对 $\omega=(\operatorname{Im}\tau)\omega_K$
  有解 $F=(\operatorname{Im}\tau)\omega_J$、$\widehat N=I$；$\omega_J$ 恰当，故 $F=\dd\alpha$ 是
  平凡线丛上的曲率，$\alpha=\frac{\operatorname{Im}\tau}{2\pi}\int_C|\dd^2z|
  \Tr(\varphi_z\delta A_{\bar z}+\varphi_{\bar z}\delta A_z)$ \pA\anchor{src p.186}。一般解
  $(\operatorname{Im}\tau)(a\omega_J+b\omega_I)$、$a^2+b^2=1$ 在 $b\neq0$ 时不恰当、周期不整，
  不能定义膜；因此 $B_{c.c.}$ 唯一（$G$ 单连通时），否则唯一到扭一个平坦线丛。它是
  $(A,B,A)$：在 $I$ 与 $K$ 中是 A-膜（$(\omega_I^{-1}\omega_J)^2=-1$），在 $J$ 中是 B-膜
  （$F\propto\omega_J$ 为 $(1,1)$ 型）\pA\anchor{src p.186}。
\item \textbf{结论}：无世界面瞬子修正——边界条件
  $G^{IJ}\partial_1\Phi_J=-\ii F^I{}_J\partial_0\Phi_J$ 与 $K$ 中全纯性联立给出
  $(\omega-\ii F)\partial_0\Phi=0$，$\omega-\ii F$ 非退化故边界上 $\Phi$ 为常数，再由解析性
  整体为常数 \pA\anchor{src p.187}。于是开弦代数只需微扰计算：作用量
  $S=-(\operatorname{Im}\tau)\int_\Sigma\Phi^*(\cos q\,\omega_J+\sin q\,\omega_I+\ii\omega_K)$
  在 $\Psi=0$ 时写成 $\omega-\ii F=-\ii\operatorname{Im}\tau\,\Omega_I$ 的势；限制到
  $T^*V\subset\MH$（$V\subset\MC$ 小开集）后，局部的乘法由 Poisson 括号
  $[\kappa,f]=-\ii(\operatorname{Im}\tau)^{-1}\kappa(f)$ 给出，而 $\hat p_\alpha=
  -\frac{\ii}{\operatorname{Im}\tau}\frac{\partial}{\partial q^\alpha}$，即全纯微分算子
  \pA\anchor{src p.191}\anchor{src p.192}。全局：粘接函数 $s^{(ij)}$ 的三重上链
  $c^{(ijk)}=s^{(ij)}+s^{(jk)}+s^{(ki)}\in H^2(\MC,\CC)$；若 $c\in2\pi\ZZ$ 得到线丛 $L$ 上的
  $\Dmod_L$，否则得到复幂 $\Dmod_{L^\gamma}$ \pA\anchor{src p.192}\anchor{src p.193}。
  时间反演 $T$（$t=1$ 处的精确对称，$(\mathbf A,\varphi)\mapsto(T^*\mathbf A,-T^*\varphi)$）把
  开弦代数映到 opposite；$\Dmod_L$ 的 opposite 是 $\Dmod_{L^{-1}\otimes K_M}$，自对偶要求
  $L^2=K_M$，即 $L=K_M^{1/2}$（$\MC$ 是 spin 流形，且平坦扭不改变层）\pA\anchor{src p.193}%
  \anchor{src p.195}。整体截面：经典时是 $\MH$ 上的全纯函数，即可积系统的交换 Hamiltonians；
  量子化即把 $w=\sum f^{\alpha\beta\gamma}p_\alpha p_\beta p_\gamma+\cdots$ 提升为符号等于 $w$
  的微分算子 $\hat w$，存在性对二阶是 Hitchin $[141]$、任意阶是 Beilinson--Drinfeld $[5]$，
  而 $T$ 给出存在性与交换性的简洁论证 \pA\anchor{src p.195}\anchor{src p.196}。
\item \textbf{出口}：\S11.2 把每个 $K$ 中的 A-膜变成 $\Dmod_{K^{1/2}}$-模；\S12.4 用四维边界条件
  整体定义 c.c.\ 膜（含奇异点与 stack）。
\item \textbf{锚点索引}：p.185--186 (11.1)--(11.3) 与 $(A,B,A)$ 判定；p.187 (11.6)--(11.7) 无瞬子；
  p.188--191 开弦代数与非交换形变；p.192--193 (11.16) 粘接与 $\Dmod_L$；p.193--195 时间反演、
  $K_{\MC}^{1/2}$、量子化；p.195--196 $[141,5]$ 与交换算子。
\end{enumerate}
\textbf{可信度}：A——余各向同性条件、无瞬子定理、局部开弦代数、时间反演钉死 $K^{1/2}$ 都是
原文给出的完整推导。B——全局层化只对 $T^*V$ 型开集成立，作者明说完全层化在精确理论中不现实
（Moyal 乘积收敛问题）；把 $\Dmod_{L}$ 的粘接当成层是微扰意义下的论断 \pB\anchor{src p.190}。
C——"Hitchin 可积系统的量子化"这一读法与 $[141,5]$ 的分工是作者的框架性判断，本节只转述
\pC\anchor{src p.196}。
\end{kcard}

\begin{kcard}{\S11.2 每个 A-膜给出一个 D-模（pp.196--201）}
\begin{enumerate}[label=\textbf{\arabic*}]
\item \textbf{这一节的任务}：把 \S11.1 的代数结论升级为函子式的陈述：任意 $K$ 中的 A-膜 $B'$
  都给出 $\Dmod_{K^{1/2}}$-模；给出两个关键例子（$\{\varphi=0\}$ 膜、type F 膜）和 $U(1)$ 的
  显式计算。
\item \textbf{输入}：\S11.1；\S10.4 用过的 $(B,A,A)$ 膜；$[142,143]$ 的开闭弦 TFT 公理；
  $[144]$ 的余各向同性膜开弦量子化。
\item \textbf{关键构造}：一般事实——任意两膜 $B,B'$ 的 $(B,B')$ 弦构成 $(B,B)$ 弦代数上的模
  （\S6.3 图 6 的切割论证），层化后即层与模层；故取 $B=B_{c.c.}$ 就把每个 $K$ 中的 A-膜 $B'$
  变成 $\Dmod_{K^{1/2}}$-模 \pA\anchor{src p.196}。$(B_{c.c.},B')$ 弦无瞬子修正（论证同 \S11.1：
  瞬子在边界上为常数，由解析性整体为常数），因此只需微扰论。例 1：$B'$ 取支撑在
  $\{\varphi=0\}\cong\MC$ 上的平坦 $(B,A,A)$ 膜；法向费米子无零模，切向零模量子化给 $\MC$ 的
  旋量 $K_{\MC}^{1/2}\otimes(\oplus_j\Omega^{0,j})$，$Q$ 化为 $\bar\partial$，故物理态就是
  $K_{\MC}^{1/2}$ 的全纯截面——这正是 $\Dmod_{K^{1/2}}$ 的定义模，反过来解释了 c.c.\ 弦代数
  为什么是它 \pA\anchor{src p.197}\anchor{src p.198}。（这里 $B'$ 在复结构 $I$ 中必须是复子流形，
  否则 $Q$ 不能化成 $\bar\partial$。）例 2：$B_F$ 是 type F 膜（Hitchin 纤维 $F$ $+$ 平坦线丛），
  它是 $(A,B,A)$ 型，故给 $\Dmod_{K^{1/2}}$-模；再加上它是 zerobrane 的 S-对偶像，就同时满足
  几何 Langlands 的两条要求：magnetic eigenbrane 与 $D$-模 \pA\anchor{src p.198}。
  $U(1)$ 显式：$\MH=\mathrm{Jac}(C)\times B$（$B=H^0(C,\Omega^1)$），$F_p\simeq\mathrm{Jac}(C)$，
  态是 $T_p=L\otimes N^{-1}|_{F_p}$ 的全纯截面；作用量 (11.20)--(11.22) 给出 $\varphi_z$ 的协变
  $\bar\partial$ 算子，即 $\mathrm{Jac}(C)$ 上平凡线丛的平坦联络，故 $D$-模对应平坦
  $\CC^*$-联络 \pA\anchor{src p.199}\anchor{src p.200}。S-对偶把这族膜对应到 $\MH(U(1),C)$ 上的
  zerobranes；而 "$\CC$ 上平坦 $\CC^*$-联络 $\leftrightarrow$ $\mathrm{Jac}(C)$ 上平坦
  $\CC^*$-联络" 就是 Poincaré 对偶 $H_1(C,\ZZ)\cong H_1(\mathrm{Jac}(C),\ZZ)$
  \pA\anchor{src p.200}\anchor{src p.201}。
\item \textbf{出口}：\S11.3 把 $\Psi=0$ 的构造推广到任意 $\Psi$；\S12 把膜整体提升回四维规范
  理论，从而 "每个 A-膜 $\Rightarrow$ $D$-模" 才在奇异点与 stack 上成立。
\item \textbf{锚点索引}：p.196 开闭弦公理与无瞬子；p.197--198 旋量与 $\bar\partial$、c.c.\ 层
  的反向解释；p.198 type F 的两条性质；p.199--200 (11.17)--(11.25) $U(1)$ 显式；p.200--201 平坦
  $\CC^*$-联络与 Poincaré 对偶。
\end{enumerate}
\textbf{可信度}：A——模结构、零模量子化、$U(1)$ 显式计算都是原文完整推导。B——"每个 A-膜给出
一个 $D$-模"是启发式陈述：层化只在 $T^*V$ 上、精确理论中未证明，作者用 "we claim/we argue"
措辞，并把整件事当作 $[41]$ 的 A-膜 $\leftrightarrow$ 非交换 B-膜这一特例
\pB\anchor{src p.196}。C——把它等同于几何 Langlands 的右端需要 $Bun_G$ 版本，本节只到
$\MC$；作者自己把这一步推给 \S12 \pC\anchor{src p.184}。
\end{kcard}

\begin{figure}[htbp]
\centering
\begin{tikzpicture}[font=\small,
  box/.style={draw=kwblue!70, fill=kwblue!4, rounded corners=1.5pt, align=center,
              inner sep=4pt, text width=3.2cm, font=\scriptsize}]
\node[box] (brane) at (0,0) {A-膜 $B'$ 于 $\MH$\\（复结构 $K$）};
\node[box] (alg) at (5.7,1.85) {c.c.\ 膜的开弦代数\\$\Dmod_{K_{\MC}^{1/2}}$};
\node[box] (mod) at (5.7,0) {$B'$ 上的开弦态\\$=$ 该代数的模};
\node[box] (dm) at (11.3,0) {$\MC$ 上 $D$-模\\（$\Psi=h$：普通 $D$-模）};
\node[box, fill=kwred!5, draw=kwred!60] (cc) at (0,1.85) {c.c.\ 膜 $B_{c.c.}$\\$(A,B,A)$，space-filling};
\draw[-{Latex[length=1.8mm]}] (cc) -- (alg) node[midway,above,font=\scriptsize]{$(B,B)$ 弦};
\draw[-{Latex[length=1.8mm]}] (brane) -- (mod) node[midway,below,font=\scriptsize]{$(B,B')$ 弦};
\draw[-{Latex[length=1.8mm]}] (alg) -- (mod);
\draw[-{Latex[length=1.8mm]}] (mod) -- (dm) node[midway,above,font=\scriptsize]{层化 / stack};
\draw[-{Latex[length=1.8mm]}, dashed, kwgreen!70!black]
  (dm.north) to[bend left=12] node[above,font=\scriptsize,text=kwgreen!70!black]{$\Psi\neq0$：扭 $\mathcal L^{\Psi-h}$} (alg.north east);
\end{tikzpicture}
\caption{\S11 的桥：c.c.\ 膜的 $(B,B)$ 弦给出微分算子层，任何 A-膜是它的模，于是得到 $\MC$ 上的
$D$-模；$\Psi=h$ 时层退化为普通微分算子 \pA\anchor{src p.196}\anchor{src p.206}。}
\label{fig:geom-dmod-bridge}
\end{figure}

\begin{kcard}{\S11.3 任意 $\Psi$：twisted D-模与 $\lambda$-连接（pp.201--206）}
\begin{enumerate}[label=\textbf{\arabic*}]
\item \textbf{这一节的任务}：把 c.c.\ 膜与 $D$-模推广到任意实 $\Psi$，得到 twisted $D$-模
  （线丛复幂上的微分算子）；并说明 $\Psi=h$ 时退化为普通 $D$-模，从而给出几何 Langlands 更
  标准的陈述。
\item \textbf{输入}：\S2.2 的 $S:\Psi\mapsto-1/n_g\Psi$；\S4.2 的 $B=-\omega_I\operatorname{Re}\tau$；
  \S6.2 的"有理 $\Psi$ 才有线算子"；$[9],[146]$。
\item \textbf{关键构造}：$\Psi\neq0$ 时 $B$-场沿 Hitchin 纤维非平坦，纤维不再是 A-膜：A-膜条件
  $F+B|_Y=0$ 要求 $[B|_Y/2\pi]$ 整，而 $B|_Y$ 是 $\omega_I$ 的倍数，$\theta\neq2\pi n$ 时无
  $U(1)$ 丛可用 \pA\anchor{src p.202}。新的 c.c.\ 膜 $B^\Psi_{c.c.}$ 由
\begin{equation}
(\omega^{-1}(F+B))^2=-1,\qquad
F=\operatorname{Im}\tau\cos q\,\omega_J,\quad
\sin q=-\frac{\operatorname{Re}\tau}{\operatorname{Im}\tau}
\label{eq:geom-cc-psi}
\end{equation}
  定义；其全纯结构是 $I(\Psi)=I\cos q-J\sin q$，全纯坐标为
  $A'_z=A_z-\ii\tan\frac q2\,\varphi_z$、$\varphi'_z=\varphi_z+\ii\tan\frac q2\,A_z$
  \pA\anchor{src p.203}。投影 $\pi_\Psi:\MH\to\MC$ 把 $(A,\varphi)$ 送到 $D'_z=D_z-\ii\tan\frac
  q2\varphi_z$ 所定全纯丛；纤维坐标是
  $D'_z=\varphi_z+\ii\tan\frac q2D_z$，即 $\lambda$-连接，$\lambda=\ii\tan\frac q2$：
\begin{equation}
\nabla(fs)=f\nabla s+\lambda\,\partial f\wedge s .
\label{eq:geom-lambda-conn}
\end{equation}
  Hitchin 方程保证 $D'_z$ 是 $\lambda$-连接，故纤维 $=$ 固定丛上 $\lambda$-连接模空间 $=$
  $H^0(C,\mathrm{End}E'\otimes K_C)$ 的仿射空间 $=$ $T^*\MC$ 的仿射丛形变
  \pA\anchor{src p.204}\anchor{src p.205}。代数结构：局部同 $\Psi=0$，全局粘接 (11.16) 被
  $\operatorname{Re}\tau$ 修正——把作用量 (11.38) 写成 $\Omega'$，其第二项是
  $\operatorname{Re}\tau$ 乘 $\MC$ 上一个非恰当的 $(1,1)$ 形式（$L$ 曲率的倍数）；摄动展开只到
  $1/(\operatorname{Im}\tau)^2$ 阶，故粘接修正对 $\Psi$ 线性；整数 $\Psi$ 处已知 $f(\Psi)=\Psi$
  （用 $\mathcal O(1)$ 上的显式模），线性性推出所有 $\Psi$ 处 $f(\Psi)=\Psi$
  \pA\anchor{src p.206}。又 $K_{\MC}\cong L^{-2h}$（$h$ 为对偶 Coxeter 数），故 c.c.\ 弦代数
  $=\Dmod_{L^{\Psi-h}}$；取 $\Psi=h$ 时就是普通微分算子 $\Dmod$ \pA\anchor{src p.206}。
\item \textbf{出口}：\S12 用四维边界条件证明这些膜整体存在（含奇异点与 stack），完成 "
  几何 Langlands $=$ $S$-对偶 $+$ $T^h$" 的表述；这也是全文末尾回扣开头 S1 的地方。
\item \textbf{锚点索引}：p.201--202 $B$-场与纤维失效；p.202--203 (11.30)--(11.33)；p.204--205
  $\lambda$-连接与 $T^*\MC$ 的仿射形变；p.205--206 粘接修正、$f(\Psi)=\Psi$；p.206
  $K_{\MC}=L^{-2h}$、$\Psi=h$、$T^hS$。
\end{enumerate}
\textbf{可信度}：A——(11.30)--(11.37) 及 $\lambda$-连接的识别是原文完整计算。B——"全局粘接修正
使 $f(\Psi)$ 线性"是二阶摄动论证（作者只算到 $1/(\operatorname{Im}\tau)^2$ 阶），且"结果对
所有 $e^2$ 成立"依赖 $\Psi$ 是唯一参数这一物理假设 \pB\anchor{src p.206}。C——"$\Psi=h$ 时
得到普通 $D$-模、这才是标准的几何 Langlands 陈述"含解释性判断，且作者指出对某些 $G$ 还要
仔细处理 $C$ 的 spin 结构依赖（引 $[5,26]$）\pC\anchor{src p.206}。
\end{kcard}

% ======================================================================
\subsection{\S12 从规范理论构造膜（pp.207--221）}

\S12 是全文的收尾，性质与前面各节不同：它不引入新的数学对象，也不给出新的对应，而是把
\S6--\S11 用到的每一个膜\emph{翻译回四维规范理论}。为什么必须做这件事？因为几何 Langlands 的
Hecke 算子作用于 $Bun_G(C)$——所有 $G$-丛的 stack，而不只是稳定丛的模空间 $\MC$；规范理论
中的对应物是"所有联络的空间"。只有把膜在四维中定义，才能说它们是在整个 stack 上定义的
\pA\anchor{src p.184}。作者在开头就把这一点说清楚：\S12 的目标就是证明"重要的膜确实可以在
四维规范理论中定义" \pA\anchor{src p.207}。

因此\S12 读起来像一张\textbf{检查清单}：对四类膜（$(B,B,B)$ 的 zerobrane、$(B,A,A)$ 的
Hitchin 纤维、$(A,B,A)$ 的 $W'$ 与 $W''$、以及 c.c.\ 膜），逐一执行同一套四步法——
(i) 选一个保超对称的对合 $T$，取费米子边界条件 $T\lambda|=\pm\lambda|$；
(ii) 要求超对称变分 $\delta\lambda=\frac12\Gamma^{IJ}F_{IJ}\epsilon$ 保持该条件，推出玻色子
边界条件；(iii) 验证超流法向分量消失 $\epsilon J_1=0$；(iv) 检查存在超对称态（$F$ 在适当子空间
自对偶）。每一步都是线性条件，可逐条核对——这正是"检查清单"的含义：它的价值在于验证整条链
自洽，而不在于产出新结论。作者也坦白：四维实现不唯一，不同边界条件在光滑点低能等价，在奇异点
可能给出不同的膜 \pA\anchor{src p.207}\anchor{src p.208}。

\begin{kcard}{\S12.1 边界条件的一般性质（pp.208--210）}
\begin{enumerate}[label=\textbf{\arabic*}]
\item \textbf{这一节的任务}：建立"边界条件 $=$ 边界数据的极大各向同性子空间"的语言，写出玻色子
  与费米子边界项、超流消失条件，为后面四步法提供验证工具。
\item \textbf{输入}：\S2.1 的场与超荷记号；\S12 开头的四维设定 $M=\RR^2_+\times C$；$N=4$ SYM 的
  拉氏量。
\item \textbf{关键构造}：最小 Yang--Mills 作用量变分的边界项为
\begin{equation}
(\delta I_B)_{\text{bdry}}=\frac{2}{e^2}\int_{\partial M}\sqrt g\,\Tr\Big(g^{IJ}\sum_{I,J\neq1}
\delta A_I F_{1J}\Big),
\label{eq:geom-bdry-term}
\end{equation}
  边界条件必须是一组使该式为零的极大线性条件（"双曲边界条件"）：$\delta A_I=0$ 即 Dirichlet，
  $F_{1I}=0$ 即规范化 Neumann \pA\anchor{src p.209}。费米子边界项为
  $\delta I_F=-\frac{\ii}{e^2}\int_{\partial M}\sqrt g\,\Tr\bar\lambda\Gamma^1\delta\lambda$，
  故只需令一半分量的边界值为零 \pA\anchor{src p.209}。超对称相容的判据是超流法向分量消失：
  $J_I=\frac12\Gamma^{KL}F_{KL}\Gamma_I\lambda$，条件 $\epsilon J_1=0$
  \pA\anchor{src p.210}。实用流程：先猜费米子条件，由 $\delta\lambda=\frac12\Gamma^{IJ}F_{IJ}
  \epsilon$ 推出玻色子条件，再验证 $J_1=0$；此外 Dirichlet 条件一般破坏超对称，需要靠 $F$ 的
  自对偶条件救回存在超对称态的可能 \pA\anchor{src p.210}。
\item \textbf{结论}：四维边界条件的合法性由三个可检验条件组成：边界项消失、$J_1=0$、存在
  超对称态。\S12.2--\S12.4 对每类膜逐一验证。
\item \textbf{出口}：\S12.2 从 $(B,B,B)$ 开始；\S12.3--\S12.4 用 S-对偶得到其余类型。
\item \textbf{锚点索引}：p.208--209 (12.2)--(12.5) 玻色/费米边界项与双曲条件；p.210
  (12.6)--(12.8) 超流与 $\delta\lambda$；p.210 超对称态的存在条件。
\end{enumerate}
\textbf{可信度}：A——边界项、超流、变分都是标准场论计算，原文完整。B——"Dirichlet 条件一般破坏
超对称、需自对偶条件补救"是经验性概括，作者没有给一般证明 \pB\anchor{src p.210}。
\end{kcard}

\begin{kcard}{\S12.2 $(B,B,B)$ 膜：zerobrane 的两种四维实现（pp.211--214）}
\begin{enumerate}[label=\textbf{\arabic*}]
\item \textbf{这一节的任务}：用四步法构造 $(B,B,B)$ 膜的边界条件，认出它就是 $\MH$ 上的
  zerobrane；并给出第二种实现，指出两者在奇异点可能不同。
\item \textbf{输入}：\S12.1 的四步法；$\Gamma^{2367}\epsilon=\epsilon$ 的紧化条件。
\item \textbf{关键构造}：取 $T=\Gamma^{04}$、费米子条件 $T\lambda|=\lambda|$。要求
  $(T-1)\delta\lambda=0$ 给出 $\sum_{I\neq0,4}(\Gamma^{0I}F_{4I}-\Gamma^{4I}F_{0I})\epsilon=0$，
  配合 $\Gamma^0\Gamma^4\epsilon=\epsilon$ 得
\begin{equation}
F_{0I}+F_{4I}=0,\qquad I\neq0,4 .
\label{eq:geom-bbb-boson}
\end{equation}
  玻色子边界条件：$A_0+A_4=0$（零方向，可同时要求其法向导数为零），$A_I$（$I\neq0,1,4$）取
  Dirichlet \pA\anchor{src p.211}\anchor{src p.212}。验证超流条件等价于 \eqref{eq:geom-bbb-boson}；
  超对称态要求 $\sum_{I,J\neq0,1,4}\Gamma^{IJ}F_{IJ}\epsilon=0$，只需 $F$ 在
  $(2,3,6,7)$ 子空间自对偶，即 Hitchin 方程
\begin{equation}
F_{23}-F_{67}=0,\qquad F_{27}-F_{36}=0,\qquad F_{26}-F_{73}=0 .
\label{eq:geom-hitchin-boundary}
\end{equation}
  故边界值给出 $\MH$ 的点 $p$；$p$ 对应不可约 Higgs 丛时 $F_{IJ}|=0$（$I$ 或 $J$ 不在
  $\{2,3,6,7\}$ 中）迫使 $A_5=A_8=A_9=0$，边界条件唯一 \pA\anchor{src p.213}。第二种构造：
  $\widetilde T=T\Gamma^{2367}=\Gamma^{023467}$，取 $\widetilde T\lambda|=-\lambda|$，得到
  (12.17)--(12.19)（$1589$ 平面的反自对偶条件），同样给 $p$ 处的 zerobrane；差别是在平坦
  时空中它保留八个超荷（前一种只保留四个），且 $A_4,A_5,A_8,A_9$ 的条件不同，故在奇异点
  （稳定化子正秩）可能不等价 \pA\anchor{src p.213}\anchor{src p.214}。第三种
  $\widetilde T\lambda|=\lambda|$ 给 space-filling 膜 $+$ 平坦规范场，原文从略
  \pA\anchor{src p.215}。
\item \textbf{结论}：zerobrane 有四维实现，且不唯一；光滑点的低能等价性可以接受，奇异点是否
  相同是开放问题。这是 \S12 检查清单的典型输出：实现存在、唯一性只是"光滑点意义下"的。
\item \textbf{出口}：\S12.3 用 S-对偶把 $(B,B,B)$ 变成 $(B,A,A)$。
\item \textbf{锚点索引}：p.211--212 (12.9)--(12.12)；p.212--213 (12.13)--(12.16) 与 Hitchin 方程；
  p.213 第二构造 $\widetilde T$；p.214--215 (12.17)--(12.19)；p.215 $A_4,A_5,A_8,A_9$ 与
  奇异点。
\end{enumerate}
\textbf{可信度}：A——(12.9)--(12.16) 与超流验证是原文完整计算。B——"两种构造在光滑点低能等价"
是原文用 "should be equivalent" 表达的推断，未证明 \pB\anchor{src p.214}。C——"在奇异点给出
不同的膜"是作者的猜测与展望 \pC\anchor{src p.214}。
\end{kcard}

\begin{kcard}{\S12.3 $(B,A,A)$ 膜：Hitchin 纤维（pp.215--217）}
\begin{enumerate}[label=\textbf{\arabic*}]
\item \textbf{这一节的任务}：对 $(B,B,B)$ 的边界条件作 S-对偶，得到 $(B,A,A)$ 膜的四维实现，
  并认出它就是 Hitchin 纤维上的膜；注意经典分析给出的方程与 Hitchin 方程差一项量子修正。
\item \textbf{输入}：\S12.2；$S$-对偶在超荷上的作用 (2.25)；\S10.5 的"支撑 Hitchin 纤维"论证。
\item \textbf{关键构造}：$S$ 把超荷按
  $\epsilon_\ell\mapsto\epsilon_\ell\ee^{-\ii\pi/4}$、$\epsilon_r\mapsto\epsilon_r
  \ee^{\ii\pi/4}$ 变换，等价地 $\epsilon\mapsto\epsilon'=\frac{1+\hat\Gamma}{\sqrt2}\epsilon$，
  其中 $\hat\Gamma=\Gamma^0\Gamma^1\Gamma^2\Gamma^3$ 测四维手征；$T$ 相应变成
  $T'=\frac{1+\hat\Gamma}{\sqrt2}T\frac{1-\hat\Gamma}{\sqrt2}$；$T=\Gamma^0\Gamma^4$ 时
  $T'=\Gamma^1\Gamma^2\Gamma^3\Gamma^4$ \pA\anchor{src p.215}。取 $T'\lambda|=-\lambda|$，
  超对称要求 $\big(\sum_{I,J\in\{1,2,3,4\}}+\sum_{I,J\notin\{1,2,3,4\}}\big)
  \Gamma^{IJ}F_{IJ}\epsilon=0$；这给出 $F_{0I}|=0$（$I>4$），$A_5=A_8=A_9=0$，以及
  $[\varphi_2,\varphi_3]=F_{67}=0$（经典）；$\varphi_z$ 的特征多项式取 Dirichlet 固定；
  $A_2,A_3,A_4$ 满足修正 Neumann 条件 (12.25) \pA\anchor{src p.216}。再要求
  $\Gamma^{2367}\epsilon=\epsilon$ 的超荷不变，得到 $F_{26}-F_{73}=F_{27}-F_{36}=0$，即
  $D_z\varphi_z=0$（复 Hitchin 方程）；于是膜支撑在 $(\varphi_z$ 特征多项式 $)$ 的纤维上，
  即 Hitchin 纤维——$(B,A,A)$ 膜的原型 \pA\anchor{src p.216}。注意经典分析给出
  $[\varphi_2,\varphi_3]=0$，而 $\MH$ 的定义是 $[\varphi_2,\varphi_3]-F_{23}=0$；原文认为
  这项差别是 IR 重整化群的量子修正 \pC\anchor{src p.216}。第二构造
  $T''=-T'\Gamma^{2367}=-\Gamma^{1467}$ 给出同样支撑在纤维上的膜，但时间反演迫使 Chan-Paton
  线丛只能是平凡或二阶平坦的；作者希望修改构造使线丛可变 \pA\anchor{src p.217}。
\item \textbf{结论}：$(B,A,A)$ 膜的 4d 实现 = 固定 $\varphi_z$ 特征多项式、让 $A_z$ 涨落，
  即 Hitchin 纤维；这是 \S10.5 中"支撑在 Hitchin 纤维上的膜不被 't Hooft 算子移出"的四维来源。
\item \textbf{出口}：\S12.4 构造 $(A,B,A)$ 与 c.c.\ 膜，完成清单。
\item \textbf{锚点索引}：p.215--216 (12.20)--(12.24)；p.216 $A_I$ 的 Dirichlet 条件与 (12.25)；
  p.216 (12.26) 复 Hitchin 与纤维支撑；p.216 量子修正；p.217 第二构造与线丛限制。
\end{enumerate}
\textbf{可信度}：A——(12.20)--(12.26) 的推导完整。B——"膜支撑在 Hitchin 纤维上"原文用
"Presumably"，是在上述线性条件读出的 \pB\anchor{src p.216}。C——$[\varphi_2,\varphi_3]=0$ 到
Hitchin 方程的量子修正是明确的猜测 \pC\anchor{src p.216}；"时间反演迫使线丛平凡/二阶"也是
作者的论断（未给完整证明）\pC\anchor{src p.217}。
\end{kcard}

\begin{kcard}{\S12.4 $(A,B,A)$ 膜与 c.c.\ 膜（pp.217--221）}
\begin{enumerate}[label=\textbf{\arabic*}]
\item \textbf{这一节的任务}：构造 $(A,B,A)$ 膜的四维边界条件，特别是 c.c.\ 膜的；并用边界算子
  分类证明 \S11.1 中 c.c.\ 弦代数的量子不变量（Q-上同调）在四维中有依据——这是 \S12 与
  \S11 的接口。
\item \textbf{输入}：\S12.1 的四步法；\S11.1 的 c.c.\ 膜与边界算子；S-对偶。
\item \textbf{关键构造}：取 $T'=-\Gamma^{15}\Gamma^{2-7}\Gamma^{3+6}$，它对 $x^2$--$x^3$ 与
  $x^6$--$x^7$ 平面同时旋转不变，故只依赖 $C$ 的复结构 \pA\anchor{src p.217}。
  $T'\lambda|=\lambda|$ 时 $A_2-\varphi_3$、$A_3+\varphi_2$ 有指定边界值，$A_2+\varphi_3$、
  $A_3-\varphi_2$ 自由涨落；故 $W'=\{A_z=A_z+\ii\varphi_z$ 固定$\}$，它在 $J$ 中是复子流形、
  对 $\Omega_J$ 拉格朗日，是支撑 $W'$ 的 $(A,B,A)$ 膜，$W'$ 参数化"固定全纯结构的平坦
  $G_{\CC}$ 丛" \pA\anchor{src p.218}。$T'\lambda|=-\lambda|$ 则交换 $A_z$ 与 $A_{\bar z}$ 的
  角色，得到 $W''$（固定反全纯结构），同样 $(A,B,A)$ \pA\anchor{src p.219}。对 $W'$ 作 S-对偶：
  $\widehat T=\frac{1+\hat\Gamma}{\sqrt2}T'\frac{1-\hat\Gamma}{\sqrt2}
  =-\frac{\Gamma^{15}}{\sqrt2}\big[(\Gamma^{26}+\Gamma^{37})
  -\hat\Gamma(\Gamma^{23}+\Gamma^{67})\big]$；$\widehat T\lambda|=-\lambda|$ 给出一组修正
  Neumann 条件 (12.31) 与 $F_{0J}=0$（$J=4,8,9$）\pA\anchor{src p.219}。丢掉大质量场
  $A_4,A_8,A_9$ 并取 $A_1=0$ 后，边界条件化为
\begin{equation}
D_0\varphi_2+\partial_1A_2=0,\quad D_0\varphi_3+\partial_1A_3=0,\quad
F_{02}-\partial_1\varphi_2=0,\quad F_{03}-\partial_1\varphi_3=0 .
\label{eq:geom-aba-bc}
\end{equation}
  与 $\sigma$-模型中的混合边界条件 $f_{MN}\partial_0X^N+g_{MN}\partial_1X^N=0$ 比较，得
\begin{equation}
f=\frac{1}{2\pi}\int_C|\dd^2z|\,\Tr\big(\delta\varphi_z\wedge\delta A_{\bar z}
+\delta\varphi_{\bar z}\wedge\delta A_z\big)=\omega_J ,
\label{eq:geom-cc-curvature}
\end{equation}
  即 c.c.\ 膜的目标空间规范场——space-filling $(A,B,A)$ 膜正在四维中被实现
  \pA\anchor{src p.220}。最后：边界算子的分类——维数与自旋都为 $n$ 的 Q-不变算子只能由
  $\varphi_z,\psi_z$ 构成（(12.35) $\psi_z|=\tilde\psi_z|$），时间反演 $T$ 给其本征值
  $(-1)^{n-k}$，故 Q 不能把 $k\to k+1$ 而不改变自旋，于是 Q-上同调 $=$ 规范不变函数
  $F(\varphi_z,\psi_z)$ 的经典结果 \pA\anchor{src p.220}\anchor{src p.221}。这正是 \S11.1
  中"把整体全纯函数量子化为微分算子"所依赖的 Q-不变性与 T 论证在四维中的来源
  \pA\anchor{src p.195}。
\item \textbf{结论}：c.c.\ 膜及其它 $(A,B,A)$ 膜都在四维中有确切边界条件；\S12 至此完成
  "四类重要膜都可在四维定义"的检查清单，全文论证闭环。
\item \textbf{出口}：全文结束。回到开头总述：S5--S8 的每一环所用的膜，都能在四维规范理论中
  找到边界条件——这是 \kw{} 把几何 Langlands 翻译成量子场论的最后一块拼图。
\item \textbf{锚点索引}：p.217--218 (12.28) 与 $W'$；p.218--219 $W''$ 与两个 $(A,B,A)$ 族；
  p.219--220 (12.30)--(12.32)；p.220 (12.33)--(12.34) c.c.\ 膜的曲率；p.220--221 (12.35) 边界
  算子分类与时间反演。
\end{enumerate}
\textbf{可信度}：A——(12.28)--(12.34) 的推导与 c.c.\ 膜的识别是原文完整计算；边界算子的
Q-上同调论证也是原文完整给出的。B——\eqref{eq:geom-cc-curvature} 与 (11.2) 的归一化相差一个
2d/4d 场归一化因子，原文未解释；把 $W'$ 的整体存在性延伸到奇异点属于推断
\pB\anchor{src p.220}。C——"丢掉大质量场后剩余条件就是 c.c.\ 膜"这一步依赖对 $e^2\to0$ 的
读法，属物理判断 \pC\anchor{src p.219}。
\end{kcard}

\begin{figure}[htbp]
\centering
\begin{tikzpicture}[font=\scriptsize,
  ckstep/.style={draw=kwblue!70, fill=kwblue!4, rounded corners=1.5pt, align=center,
               inner sep=3pt, text width=2.55cm, minimum height=1.05cm}]
\node[ckstep] (s1) at (0,0) {选对合 $T$\\（保超对称）};
\node[ckstep] (s2) at (3.1,0) {猜费米子边界\\$T\lambda|=\pm\lambda|$};
\node[ckstep] (s3) at (6.2,0) {$\delta\lambda=\frac12\Gamma^{IJ}F_{IJ}\epsilon$\\$\Rightarrow$ 玻色子边界};
\node[ckstep] (s4) at (9.3,0) {验证 $\epsilon J_1=0$\\（超流法向消失）};
\node[ckstep] (s5) at (12.4,0) {检查超对称态\\（$F$ 自对偶）};
\foreach \a/\b in {s1/s2,s2/s3,s3/s4,s4/s5}
  \draw[-{Latex[length=1.6mm]}] (\a) -- (\b);
\node[draw=kwred!60, fill=kwred!5, rounded corners=1.5pt, align=center,
      inner sep=3pt, text width=2.6cm, font=\scriptsize] (out) at (12.4,-1.75)
  {得到膜型：\\$(B,B,B)$\\$(B,A,A)$\\$(A,B,A)$};
\draw[-{Latex[length=1.6mm]}, kwred!70] (s5) -- (out);
\node[font=\scriptsize, text=kwgray] at (0,-1.15) {四步法（每步线性、可逐条核对）};
\end{tikzpicture}
\caption{\S12 的检查清单结构：同一套四步法在 \S12.2--\S12.4 中分别作用于三类膜。它不产生新
结论，只验证 \S6--\S11 的每个膜都能在四维规范理论中定义 \pA\anchor{src p.207}\anchor{src p.208}。}
\label{fig:geom-checklist}
\end{figure}

% ======================================================================
\subsection*{本节锚点索引}

\begingroup\small
\begin{longtable}{@{}p{2.2cm}p{1.5cm}p{4.9cm}p{6.4cm}@{}}
\toprule
\textbf{卡片} & \textbf{印刷页} & \textbf{原文锚点（定位用）} & \textbf{核心断言} \\
\midrule
\endhead
\S6.1 Wilson 线 & p.79--p.80 & p.79 (6.1)；p.80 (6.2)--(6.5) & Wilson 算子 $=$ 沿 $S=\gamma\times p$ 的平行移动；$t=\pm\ii$ 时拓扑。\\
\S6.2 't Hooft 算子 & p.81--p.85 & p.81 (6.7)--(6.9)；p.84 (6.14)--(6.16)；p.85 GNO & 't Hooft 算子按 $\mathrm{Rep}(\Ldual)$ 分类；有理 $\Psi$ 才有混合算子。\\
\S6.3 线算子 $=$ 膜 & p.89--p.92 & p.90 三类曲线；p.91 折叠图 5；p.92 (6.17)--(6.19) & 线算子与 $(Y,Z)$-膜等价；拼接代数 $=$ 开弦代数。\\
\S6.4 特征膜 & p.93--p.102 & p.95 (6.22)；p.97 (6.25)；p.98 (6.27)；p.99 (6.29) & eigenbrane $=$ 带平坦丛的膜；$S$ 交换电/磁；$\mu$OPE 为 "we expect"。\\
\S7.1 中心与万有丛 & p.103--p.109 & p.105 中心表；p.106 (7.4)--(7.7)；p.108 (7.12) & 中心 $Z$ 自对偶；universal bundle 障碍 $\zeta\in H^2(M,Z)$；twisted 丛。\\
\S7.2 flux 的 2d 含义 & p.110--p.116 & p.111 (7.18)--(7.22)；p.112 B-field；p.114 HT & $e_0,m_0,e_1,m_1$ 的分解；$S$ 交换 $B$-场与拓扑类；T-对偶。\\
\S8.1 Wilson 作用于膜 & p.116--p.120 & p.117 (8.1)；p.118 (8.2)--(8.4)；p.120 (8.8)--(8.9) & Chan-Paton 丛 $U\mapsto U\otimes E_p(R)$，扭量 $=\theta_R(\zeta)$。\\
\S8.2 zerobrane & p.121--p.123 & p.121--p.122 (8.10)；p.123 GL 主张 & zerobrane $=$ electric eigenbrane；type F 膜 $=$ magnetic eigenbrane。\\
\S9.1 算子 $=$ Hecke & p.128--p.134 & p.130 (9.9)；p.131--p.132 (9.10)；p.133--134 & $T(\Lw)$ 的作用 $=L_i\mapsto L_i\otimes\mathcal O(p_0)^{m_i}$。\\
\S9.2 修改空间 & p.135--p.144 & p.136--p.137 局部模型；p.139--p.140 紧化；p.144 (9.25) & $Y(\Lw)$ 有限维、一般奇异；紧化 $=$ 加入 associated weights。\\
\S9.3 Schubert 胞腔 & p.145--p.148 & p.145--p.146 $\mathrm{Gr}_N$；p.147--p.148 轨道与闭包 & $Y(\Lw)=$ 仿射 Grassmannian 的 Schubert 胞腔。\\
\S10.1 边界条件与指标 & p.149--p.156 & p.151 两类问题；p.152 (10.1)；p.154--p.156 & 椭圆边界条件使 $Z$ 光滑部分 $=$ Hecke 空间；$\dim Z=\sum\Delta_\rho$。\\
\S10.2 单极子成泡 & p.155--p.159 & p.156--p.157 成泡；p.158 (10.12)--(10.13) & 成泡 $=$ associated weights 的加入；$Z$ 的奇点超凯勒。\\
\S10.3 Kähler 与 OPE 极限 & p.159--p.166 & p.160--p.162 (10.16)--(10.19)；p.165 (10.20) & $Z$ 的复结构 $=$ Hecke 空间；$p'\to p$ 极限展开 OPE。\\
\S10.4 't Hooft OPE & p.167--p.178 & p.169 (10.21)--(10.23)；p.172--p.174 & 系数空间 $=$ 张量积重数空间 $N_\alpha$（$[96$--$98]$）。\\
\S10.5 \mbox{扩展方程} & p.179--p.183 & p.181 (10.36)；p.182 (10.39) & 时间无关解 $=$ Higgs 丛的 Hecke 修改；保 Hitchin 纤维。\\
\S11.1 c.c.\ 膜 & p.185--p.196 & p.185--p.186 (11.1)--(11.3)；p.192 (11.16) & c.c.\ 弦代数 $=\Dmod_{K_{\MC}^{1/2}}$；时间反演钉死线丛。\\
\S11.2 膜 $\Rightarrow$ $D$-模 & p.196--p.201 & p.196--p.197 模结构；p.198 type F；p.200 $U(1)$ & 每个 A-膜是 $\Dmod_{K^{1/2}}$-模；type F 给出 GL 的 $D$-模。\\
\S11.3 twisted $D$-模 & p.201--p.206 & p.202--p.203 (11.30)--(11.33)；p.206 $L^{-2h}$ & 任意 $\Psi$：$\Dmod_{L^{\Psi-h}}$；$\Psi=h$ 时为普通 $D$-模。\\
\S12.1 边界条件一般性 & p.208--p.210 & p.209 (12.2)--(12.5)；p.210 (12.6)--(12.8) & 边界条件 $=$ 极大各向同性；四步法（$J_1=0$）。\\
\S12.2 $(B,B,B)$ & p.211--p.214 & p.212 (12.11)--(12.12)；p.213 (12.16)；p.214 (12.19) & zerobrane 的两种四维实现；光滑点等价、奇异点待定。\\
\S12.3 $(B,A,A)$ & p.215--p.217 & p.215--p.216 (12.20)--(12.26)；p.216 量子修正 & 固定 $\varphi_z$ 特征多项式 $\Rightarrow$ 支撑在 Hitchin 纤维；差项为 IR 修正。\\
\S12.4 $(A,B,A)$ 与 c.c. & p.217--p.221 & p.218 (12.28)；p.220 (12.34)；p.221 (12.35) & c.c.\ 膜的 4d 边界条件；边界算子 Q-上同调补上 \S11.1。\\
\bottomrule
\end{longtable}
\endgroup

\noindent 全节共 23 张卡片（\S6.1--\S12.4）、9 张 tikz 图、208 个原文锚点，其中标注 A 级 114 处、
B 级 27 处、C 级 21 处；锚点按印刷页标注，正文在
\texttt{source/text/p\{$N+1$:03d\}.txt}。证据分级沿用前言约定：\pA{} 原文明确写出并给出
（或可复算的）推导；\pB{} 原文暗示、由我们重建或原文只给启发式论证；\pC{} 解释性判断与作者
明确标注为期待（"we expect/we believe"）的猜测。全节 A/B/C 的分布显示：算子层的
\S6--\S10 有大量 A 级严格计算，而 \S11--\S12 的"膜 $\Rightarrow$ $D$-模"与"边界条件 $\Rightarrow$
膜型"两步以 B 级物理推断为主——这正是数学读者读后半程时最需要保持警惕的两处。
